\chapter{THE STABILITY OF JETS AND CYLINDERS}
\label{ch:12}

\section{Introduction}\label{sec:107}
In this chapter we turn our attention to a different group of problems: the stability of jets and cylinders. As Lord Rayleigh has remarked in his \emph{Theory of Sound}: `a large class of phenomena, interesting not only in themselves but also in throwing light upon others yet more obscure, depend for their explanation upon the transformations undergone by a cylindrical body of liquid when slightly displaced from its equilibrium configuration and left to itself.' The classic example in this connexion is, of course, the instability of water issuing from a nozzle as a cylindrical jet. The cause of this instability, as Plateau first showed, is surface tension which makes the infinite cylinder an unstable figure of equilibrium and entails its breaking into separate pieces with a total superficial area which is less than that of the original cylinder. Specifically, Plateau showed that if the equilibrium surface $\varpi = R$ is deformed so that it becomes
\begin{equation}
\label{eq:12-1}
\varpi = R + \epsilon \cos kz,
\end{equation}
where $\epsilon$ is a small constant, $z$ is measured along the axis, and $\varpi$ is the horizontal distance from the axis, the deformation is stable or unstable according as $kR$ is greater than or less than unity, i.e.\ according as the wavelength of the varicosity\footnote{In introducing this terminology to describe the deformation (1), Lord Rayleigh has said: `In speaking of this subject, I have often been embarrassed for want of an appropriate word to describe the condition in question. But a few days ago, during a biological discussion, I found that there is a recognized if not a very pleasant word. The cylindrical jet may be said to become varicose and the varicosity goes on increasing until eventually it leads to absolute disruption.' In recent times, `sausage instability' has been used to describe the same condition; but this is also not a very `pleasant' description, and varicose instability would seem preferable.} is less than or greater than the circumference of the cylinder. Plateau concluded from this result that a cylindrical jet will break-up into pieces of length $2\pi R$. This conclusion, as Rayleigh showed, is incorrect: for, to ascertain the manner in which an unstable system will tend to break-up, one must determine the mode of \emph{maximum instability}. The principle here is the following. If a system is characterized by a number of unstable modes symbolized by $k_1, k_2, \ldots, k_n$, the amplitudes of which increase like $e^{\sigma_1 t}, e^{\sigma_2 t}, \ldots, e^{\sigma_n t}$ (where all the $\sigma$'s have positive real parts), then the mode of maximum instability is that one which belongs to the $\sigma$ with the largest real part. It is by this mode the system will increasingly tend to break-up when an arbitrary initial disturbance is reduced to evanescence: for, if $k_0$ denotes the mode of maximum instability, then its amplitude, relative to another mode $k_s$, after a certain time $t$, will be $\exp[\mathrm{re}(\sigma_0 - \sigma_s)t]$; and this tends to infinity as $t$ tends to infinity.

Two principal problems with which we shall be concerned in this Chapter (\S\S~108--12) are the \emph{gravitational} and the \emph{capillary} instability of an infinite liquid cylinder. The two problems will be considered in the frameworks of hydrodynamics and hydromagnetics, when dissipative mechanisms are not present and when they are present. In \S\S~113 and 114 we shall consider the effect of fluid motions on the hydromagnetic stability of certain cylindrical systems; and in \S~115 we shall consider the stability of the so-called pinch configuration.

\section{The gravitational instability of an infinite cylinder}\label{sec:108}
Consider a uniform infinite cylinder of an incompressible, inviscid fluid. The equations governing small departures from equilibrium are
\begin{equation}
\label{eq:12-2}
\frac{\partial \mathbf{u}}{\partial t} = -\operatorname{grad} \Pi,
\end{equation}
\begin{equation}
\label{eq:12-3}
\operatorname{div} \mathbf{u} = 0, \quad \nabla^2 \mathfrak{B} = 0,
\end{equation}
and
\begin{equation}
\label{eq:12-4}
\Pi = -\delta\mathfrak{B} + \delta p / \rho,
\end{equation}
where the various symbols have their usual meanings.

Since we are considering departures from an equilibrium right-cylindrical shape of an incompressible fluid, a normal mode can be expressed uniquely in terms of the deformed surface. Suppose then that the deformed surface is described by the equation
\begin{equation}
\label{eq:12-5}
\varpi = R + \epsilon e^{i(kz + m\varphi)},
\end{equation}
where $R$ is the radius of the unperturbed cylinder, $\epsilon$ is some function of the time (which we shall presently specify to be of the form $\epsilon_0 e^{\sigma t}$), $k$ is the wave number of the disturbance in the $z$-direction, and $m$ is an integer (positive, zero, or negative).

For deformations of the boundary described by equation~\eqref{eq:12-5}, the change in the gravitational potential can be obtained as follows. The external and the internal gravitational potentials satisfy the Laplace and the Poisson equations
\begin{equation}
\label{eq:12-6}
\nabla^2 \mathfrak{B}_{\mathrm{ext}} = 0 \quad \text{and} \quad \nabla^2 \mathfrak{B}_{\mathrm{int}} = -4\pi G\rho.
\end{equation}
The solutions of these equations appropriate to the deformed boundary~\eqref{eq:12-5} are
\begin{equation}
\label{eq:12-7}
\mathfrak{B}_{\mathrm{ext}} = -2\pi G\rho \varpi^2 + A K_m(k\varpi) e^{i(kz+m\varphi)} + c,
\end{equation}
and
\begin{equation}
\label{eq:12-8}
\mathfrak{B}_{\mathrm{int}} = \pi G\rho(r_1^2 - \varpi^2) + B I_m(k\varpi) e^{i(kz+m\varphi)},
\end{equation}
where $I_m$ and $K_m$ are the Bessel functions of order $m$ for a purely imaginary argument, $A$ and $B$ are constants to be determined, and $c$ is an additive constant with which we need not be further concerned. The constants $A$ and $B$ are to be determined by the condition that the potential (with a suitable fixed choice of reference) and its derivative are continuous on the boundary~\eqref{eq:12-5}. To the first order in $\epsilon$, these conditions give
\begin{equation}
\label{eq:12-9}
A K_m(x) = B I_m(x),
\end{equation}
and
\begin{equation}
\label{eq:12-10}
A K'_m(x) - B I'_m(x) = -4\pi G\rho R \epsilon,
\end{equation}
where
\begin{equation}
\label{eq:12-11}
x = kR.
\end{equation}
On solving equations~\eqref{eq:12-9} and \eqref{eq:12-10}, we find\footnote{In obtaining these solutions, use has been made of the relation $I_m(x) K'_m(x) - K_m(x) I'_m(x) = -1/x$.}
\begin{equation}
\label{eq:12-12}
A = 4\pi G\rho R_\epsilon I_m(x) \quad \text{and} \quad B = 4\pi G\rho R\epsilon K_m(x).
\end{equation}
The solution for $\delta\mathfrak{B}$ is therefore given by
\begin{equation}
\label{eq:12-13}
\delta\mathfrak{B} = 4\pi G\rho R_\epsilon K_m(x) I_m(k\varpi) e^{i(kz+m\varphi)}.
\end{equation}
Returning to equation~\eqref{eq:12-2} and taking the divergence of this equation, we obtain
\begin{equation}
\label{eq:12-14}
\nabla^2 \Pi = 0.
\end{equation}
Under the present circumstances the required solution for $\Pi$ must, therefore, be of the form
\begin{equation}
\label{eq:12-15}
\Pi = \Pi_0 I_m(k\varpi) e^{i(kz+m\varphi)},
\end{equation}
where $\Pi_0$ is a constant, unspecified for the present.

We shall now specify the dependence on time of the various quantities by assuming that
\begin{equation}
\label{eq:12-16}
\epsilon = \epsilon_0 e^{\sigma t},
\end{equation}
where $\sigma$ is a constant. Equation~\eqref{eq:12-2} then gives
\begin{equation}
\label{eq:12-17}
\mathbf{u} = -\frac{\sigma}{\Pi_0}\operatorname{grad}[I_m(k\varpi)e^{i(kz+m\varphi)}].
\end{equation}
In particular,
\begin{equation}
\label{eq:12-18}
u_\varpi = -\frac{\sigma}{\Pi_0} \frac{k}{I_m(k\varpi)} I'_m(k\varpi) e^{i(kz+m\varphi)}.
\end{equation}

The boundary conditions which must be satisfied by the solution represented by equation~\eqref{eq:12-18} are: (1) the radial component of the velocity given by~\eqref{eq:12-18} must be compatible with the assumed form of the deformed surface, namely,
\begin{equation}
\label{eq:12-19}
\dot{\varpi} = R + \epsilon_0 \sigma e^{i(kz+m\varphi)+\sigma t};
\end{equation}
and (2) the pressure must vanish on the same boundary. The first of these conditions clearly requires
\begin{equation}
\label{eq:12-20}
-\frac{\epsilon_0 \Pi_0}{\sigma} k I'_m(x) = \epsilon_0 \sigma
\end{equation}
or
\begin{equation}
\label{eq:12-21}
\sigma^2 = -\frac{\Pi_0}{\epsilon_0} k I'_m(x).
\end{equation}

Now the pressure in the unperturbed cylinder is given by
\begin{equation}
\label{eq:12-22}
p_0 = \pi G\rho^2(R^2 - \varpi^2).
\end{equation}
The vanishing of the pressure, $p_0 + \delta p$, in the perturbed cylinder, on the boundary~\eqref{eq:12-19} gives
\begin{equation}
\label{eq:12-23}
-2\pi G\rho R_\epsilon e^{i(kz+m\varphi)+\sigma t} + \left(\frac{\delta p}{\rho}\right)_R = 0;
\end{equation}
whereas according to equations~\eqref{eq:12-4}, \eqref{eq:12-13}, and \eqref{eq:12-15},
\begin{equation}
\label{eq:12-24}
-4\pi G\rho R_\epsilon K_m(x) I_m(x) e^{i(kz+m\varphi)+\sigma t} + \epsilon_0 \Pi_0 I_m(x) e^{i(kz+m\varphi)+\sigma t} = 0.
\end{equation}
Eliminating $(\delta p)_R$ from equations~\eqref{eq:12-23} and \eqref{eq:12-24}, we find
\begin{equation}
\label{eq:12-25}
4\pi G\rho R\epsilon K_m(x) I_m(x) - \frac{\Pi_0}{\epsilon_0} = -\Pi_0 I_m(x).
\end{equation}
From equations~\eqref{eq:12-21} and \eqref{eq:12-25} we now obtain
\begin{equation}
\label{eq:12-26}
\frac{\sigma^2}{4\pi G\rho} = \frac{x I'_m(x)}{I_m(x)} \left[ K_m(x) I_m(x) - \frac{1}{2} \right],
\end{equation}
where we have written $\sigma_m$ in place of $\sigma$ to distinguish the modes of different $m$'s.

From the known properties of the Bessel functions it follows that
\begin{equation}
\label{eq:12-27}
K'_m(x) I_m(x) < \tfrac{1}{2} \quad \text{for all } m \neq 0.
\end{equation}
Hence
\begin{equation}
\label{eq:12-28}
\sigma_m^2 < 0 \quad \text{for all } m \neq 0.
\end{equation}
\emph{The cylinder is, therefore, stable for all purely non-axisymmetric perturbations.}

Considering the case $m = 0$ (when the inequality~\eqref{eq:12-27} does not hold for all $x$), we can write (since $I_0(x) = J_0(ix)$),
\begin{equation}
\label{eq:12-29}
\frac{\sigma_0^2}{4\pi G\rho} = \frac{x I_1(x)}{I_0(x)} \left[ K_0(x) I_0(x) - \tfrac{1}{2} \right].
\end{equation}

The asymptotic behaviours of the different Bessel functions which appear in equation~\eqref{eq:12-29} are
\[
I_0(x) \to 1, \quad I_1(x) \to \tfrac{1}{2}x, \quad \text{and} \quad K_0(x) \to -(\gamma + \log \tfrac{1}{2}x), \quad \text{for } x \to 0,
\]
and
\begin{equation}
\label{eq:12-30}
I_0(x) \sim \frac{e^x}{\sqrt{2\pi x}}, \quad I_1(x) \sim \frac{e^x}{\sqrt{2\pi x}}, \quad \text{and} \quad K_0(x) \sim e^{-x}\sqrt{\frac{\pi}{2x}}
\end{equation}
\begin{equation}
\label{eq:12-31}
\text{for } x \to \infty.
\end{equation}
where $\gamma$ $(= 0.5772\ldots)$ is Euler's constant, and

Hence, while $I_0(x)K_0(x)$ tends to infinity logarithmically for $x \to 0$, it tends to zero as $1/(2x)$ for $x \to \infty$; moreover, $I_1(x)/I_0(x)$ is a monotone decreasing function of $x$. From these facts it follows that the equation
\begin{equation}
\label{eq:12-32}
I_0(x)K_0(x) - \tfrac{1}{2} = 0
\end{equation}
allows a single positive root. If $x_a$ denotes this root,
\[
\sigma_0^2 > 0 \quad \text{for } 0 < x < x_a \quad \text{and} \quad \sigma_0^2 < 0 \quad \text{for } x > x_a.
\]
This clearly implies that the cylinder is gravitationally unstable for all varicose deformations with wavelengths exceeding
\begin{equation}
\label{eq:12-33}
\lambda_a = 2\pi R / x_a,
\end{equation}
where $x_a$ is the root of equation~\eqref{eq:12-32}. Numerically, it is found that (see Table~LIII)
\begin{equation}
\label{eq:12-34}
x_a = 1.0668.
\end{equation}

\begin{table}[ht]
\centering
\caption{The dispersion relation for the gravitationally unstable modes of an infinite cylinder}
\label{tab:LIII}
\begin{tabular}{cccc}
\hline
$x$ & $\frac{xI_1(x)}{I_0(x)}$ & $I_0(x)K_0(x)$ & $\frac{xI_1(x)}{I_0(x)}[K_0(x)I_0(x)-\frac{1}{2}]$ \\
\hline
0 & 0 & $\infty$ & 0 \\
0.25 & 0.1253 & 1.6950 & 0.150 \\
0.50 & 0.2520 & 1.1145 & 0.155 \\
0.75 & 0.3815 & 0.8105 & 0.118 \\
0.80 & 0.4073 & 0.7640 & 0.107 \\
0.85 & 0.4333 & 0.7204 & 0.095 \\
0.90 & 0.4596 & 0.6795 & 0.083 \\
0.95 & 0.4861 & 0.6411 & 0.069 \\
1.00 & 0.5130 & 0.6050 & 0.054 \\
1.0668 & 0.5507 & 0.5000 & 0 \\
\hline
\end{tabular}
\end{table}

Table~LIII gives the value of $\sigma_0^2/(4\pi G\rho)$ for values of $x$ in the range $0 \leq x \leq 1.067$ (see also Fig.~123). It will be observed that in this range of $x$, $\sigma_0$ attains a maximum value ($\sim 0.2455\sqrt{4\pi G\rho}$) at $x = 0.580$. By arguments which have been stated in \S~107, it is by this mode of maximum instability that an infinite cylinder will break-up gravitationally when

\figurePlaceholder{123}{The dispersion relation for the gravitationally unstable modes ($x < 1.067$) of an infinite cylinder. The abscissa measures the wave number in the unit $1/R$ and the ordinate the growth rate in the unit $\sqrt{4\pi G\rho}$. The mode of maximum instability occurs for $x = 0.580$.}

the amplitude of the initial disturbance tends to zero; and the characteristic time for break-up is
\begin{equation}
\label{eq:12-35}
\tau \approx \frac{1}{0.2455\sqrt{4\pi G\rho}}.
\end{equation}

For a cylinder of radius 250 parsecs ($7.7 \times 10^{20}$~cm) and density $\rho = 2 \times 10^{-24}$~g/cm$^3$, the wavelength of the mode of maximum instability is $2.7 \times 10^3$ parsecs; and the corresponding characteristic time for break-up is $1.0 \times 10^8$ years.

\subsection*{(a) The origin of the gravitational instability and an alternative method of determining the characteristic frequencies}

The origin of the gravitational instability of an infinite cylinder can be traced directly to the sign of the change in the gravitational potential energy $\Delta\mathfrak{B}$ per unit length resulting from various deformations of different wavelengths. We shall find that $\Delta\mathfrak{B} < 0$ for precisely those modes for which the cylinder is unstable according to equations~\eqref{eq:12-28} and \eqref{eq:12-29}.

Since the potential energy per unit length of an infinite cylinder is infinite, the evaluation of $\Delta\mathfrak{B}$ requires some care. We proceed as follows.

In evaluating $\Delta\mathfrak{B}$, it is convenient to suppose that the equation of the deformed surface is given in the real form
\begin{equation}
\label{eq:12-36}
\varpi = R + \epsilon \cos kz \cos m\varphi.
\end{equation}
Suppose further that the amplitude of the deformation is increased by $\delta\epsilon$. The change $\delta\Delta\mathfrak{B}$ in the potential energy consequent to this infinitesimal increase in the amplitude of the deformation can be determined by evaluating the work done in the redistribution of the matter required to effect the change in $\epsilon$. For this evaluating this work it is necessary to specify in a quantitative manner the redistribution which does take place; and we shall now do this.

An arbitrary deformation of an incompressible fluid can be thought of as resulting from a Lagrangian displacement $\boldsymbol\xi$ applied to each point of the fluid. The incompressibility of the medium requires that $\operatorname{div}\boldsymbol\xi = 0$; and since no loss of generality is implied by supposing that the displacement is carried out irrotationally, we shall write
\begin{equation}
\label{eq:12-37}
\boldsymbol\xi = \operatorname{grad} \phi
\end{equation}
and require that
\begin{equation}
\label{eq:12-38}
\nabla^2 \phi = 0.
\end{equation}
A solution of equation~\eqref{eq:12-38} which is suitable for considering the deformation of a uniform cylinder into one whose boundary is given by~\eqref{eq:12-36} is
\begin{equation}
\label{eq:12-39}
\phi = A I_m(k\varpi) \cos kz \cos m\varphi,
\end{equation}
where $A$ is a constant. The corresponding radial component of $\boldsymbol\xi$ is
\begin{equation}
\label{eq:12-40}
\frac{\partial \phi}{\partial \varpi} = Ak I'_m(k\varpi) \cos kz \cos m\varphi.
\end{equation}
Since at $\varpi = R$, $\xi_\varpi$ must reduce to $\epsilon \cos kz \cos m\varphi$,
\begin{equation}
\label{eq:12-41}
A = \frac{\epsilon}{k I'_m(x)}.
\end{equation}
Hence,
\begin{equation}
\label{eq:12-42}
\phi = \frac{\epsilon}{k I'_m(x)} I_m(k\varpi) \cos kz \cos m\varphi
\end{equation}
\begin{equation}
\label{eq:12-43}
\boldsymbol\xi = \frac{\epsilon}{k I'_m(x)} \operatorname{grad}[I_m(k\varpi) \cos kz \cos m\varphi].
\end{equation}
Accordingly, the displacement $\delta\boldsymbol\xi$ which must be applied to each point of the fluid to increase the amplitude of the deformation by an amount $\delta\epsilon$ is given by
\begin{equation}
\label{eq:12-44}
\delta\boldsymbol\xi = \frac{\delta\epsilon}{k I'_m(x)} \operatorname{grad}[I_m(k\varpi) \cos kz \cos m\varphi].
\end{equation}

The change in the potential energy $\delta\Delta\mathfrak{B}$ per unit length involved in the additional deformation $\delta\epsilon$ can be obtained now by integrating, over the whole cylinder, the work done by the displacement $\delta\boldsymbol\xi$ in the gravitational potential [cf.\ equations~\eqref{eq:12-6} and \eqref{eq:12-13}]
\begin{equation}
\label{eq:12-45}
\mathfrak{B} = -\pi G\rho \varpi^2 + 4\pi G\rho R_\epsilon K_m(x) I_m(k\varpi) \cos kz \cos m\varphi.
\end{equation}
Thus,
\begin{equation}
\label{eq:12-46}
\delta\Delta\mathfrak{B} = -2\rho \overline{\int_0^{R+\epsilon\cos kz\cos m\varphi} (\boldsymbol\xi \cdot \operatorname{grad} \mathfrak{B})\varpi\, d\varpi},
\end{equation}
where the angular brackets signify that the quantity enclosed should be averaged over all $z$ and $\varphi$.

Substituting for $\mathfrak{B}$ and $\delta\boldsymbol\xi$ in~\eqref{eq:12-46} and neglecting terms of order $\epsilon^3$ and higher, we find,
\begin{equation}
\label{eq:12-47}
\delta\Delta\mathfrak{B} = -2\pi\rho \delta_{\epsilon,m} \left\{(-2\pi G\rho)\frac{2\rho R_\epsilon}{\sigma k I'_m(x)} \cos kz \cos m\varphi \,d\varpi\right\},
\end{equation}
where
\[
\delta_{\epsilon,m} = 1 \quad \text{if } m = 0 \quad \text{and zero otherwise.}
\]
The reduction of the first of the two integrals in~\eqref{eq:12-47} is immediate and we can write
\begin{equation}
\label{eq:12-48}
\delta\Delta\mathfrak{B} = -(1+\delta_{\epsilon,m})4\pi^2 G\rho^2 R^2 \epsilon\, \delta\epsilon \int_0^1 \left[\frac{x}{I'_m(x)} I'_m(xy) + \left(1 + \frac{m^2}{x^2 y^2}\right) I_m(xy)\right] y\, dy;
\end{equation}
and making use of the identity (which follows from Bessel's equation)
\begin{equation}
\label{eq:12-49}
\frac{d}{dy}[y I'_m(xy)] = x\left[I''_m(xy) + \left(1 + \frac{m^2}{x^2 y^2}\right) I_m(xy)\right],
\end{equation}
we obtain
\begin{equation}
\label{eq:12-50}
\delta\Delta\mathfrak{B} = -(1+\delta_{\epsilon,m}) 4\pi^2 G\rho^2 R^2 \left[\frac{x}{I'_m(x)} K'_m(x) I_m(x) - \frac{1}{2}\right] \epsilon\, \delta\epsilon;
\end{equation}
and integrating this last expression from 0 to $\epsilon$, we finally obtain
\begin{equation}
\label{eq:12-51}
\Delta\mathfrak{B} = -(1+\delta_{\epsilon,m}) 2\pi^2 G\rho^2 R^2 \left[K_m(x) I_m(x) - \frac{1}{2}\right] \epsilon^2.
\end{equation}

Comparing equation~\eqref{eq:12-51} for $\Delta\mathfrak{B}$ with equation~\eqref{eq:12-26} for $\sigma_0^2$, we observe that the criterion for stability which follows from the conditions $\Delta\mathfrak{B} > 0$ and $\sigma_m^2 < 0$ is the same. It is, of course, necessary on physical grounds that the two conditions lead to the same criterion; that they must is often referred to as the \emph{energy principle}.

The foregoing method of determining the criterion for stability by analysing the sign of the change in $\Delta\mathfrak{B}$ consequent to a particular perturbation can be extended to provide a complete analytical solution to the problem. For this purpose we suppose the amplitude $\epsilon$ of the deformation is a function of time and, considering it as a Lagrangian coordinate, we seek an equation of motion for it by defining a suitable Lagrangian function.

When $\epsilon$ is a function of time, each element of the fluid will execute motions; these can be derived from the Lagrangian displacement $\boldsymbol\xi$ by
\begin{equation}
\label{eq:12-52}
\mathbf{u} = \dot{\epsilon} \frac{\partial \boldsymbol\xi}{\partial \epsilon}.
\end{equation}

With $\boldsymbol\xi$ given by equation~\eqref{eq:12-43},
\begin{equation}
\label{eq:12-53}
\mathbf{u} = -\frac{\dot{\epsilon}}{k I'_m(x)}[\operatorname{grad}(I_m(k\varpi) \cos kz \cos m\varphi)] \frac{\partial \epsilon}{\partial \epsilon}.
\end{equation}
The kinetic energy $\mathfrak{T}$ per unit length, associated with the motions specified by~\eqref{eq:12-52}, is
\begin{equation}
\label{eq:12-54}
\mathfrak{T} = \tfrac{1}{2}(1+\delta_{\epsilon,m})\pi R^2 \rho \dot{\epsilon}^2 \left\langle \left\{ \operatorname{grad}(I_m(k\varpi) \cos kz \cos m\varphi) \right\}^2 \right\rangle \varpi\, d\varpi;
\end{equation}
or, making use of the same identity~\eqref{eq:12-49}, we obtain
\begin{equation}
\label{eq:12-55}
\mathfrak{T} = \tfrac{1}{2}(1+\delta_{\epsilon,m})\pi R^2 \rho \dot{\epsilon}^2 \frac{I_m(x)}{x I'_m(x)}.
\end{equation}
Combining this expression for the kinetic energy with the corresponding expression for the potential energy given by~\eqref{eq:12-51}, we can now define the Lagrangian function
\begin{equation}
\label{eq:12-56}
\mathfrak{L} = \tfrac{1}{2}(1+\delta_{\epsilon,m})\pi R^2 \rho \left\{ \frac{I_m(x)}{x I'_m(x)} \dot{\epsilon}^2 + 4\pi G\rho R^2 [K_m(x) I_m(x) - \tfrac{1}{2}] \epsilon^2 \right\}.
\end{equation}
The equation of motion for $\epsilon$ which follows from this Lagrangian is
\begin{equation}
\label{eq:12-57}
\frac{I_m(x)}{x I'_m(x)} \ddot{\epsilon} = 4\pi G\rho \left[K_m(x) I_m(x) - \tfrac{1}{2}\right] \epsilon.
\end{equation}
Therefore,
\begin{equation}
\label{eq:12-58}
\epsilon = \epsilon_0 e^{\sigma t},
\end{equation}
where $\sigma$ is given by the same formula~\eqref{eq:12-26}. While this method of deriving $\sigma$ is perhaps less direct analytically, it is more directly physical, since the origin of the instability is traced to the energetics of the problem.

\section{The effect of viscosity on the gravitational instability of an infinite cylinder}\label{sec:109}
We shall now consider the effect of viscosity on the gravitational instability of an infinite cylinder. Since the principal interest is in the unstable modes, we shall restrict the present discussion to axisymmetric perturbations. It will be convenient then to use the representation of axisymmetric solenoidal vectors described in Chapter~VI, \S~61~(a).

When allowance is made for viscosity, equation~\eqref{eq:12-2} of \S~108 is replaced by
\begin{equation}
\label{eq:12-59}
\frac{\partial \mathbf{u}}{\partial t} = -\operatorname{grad} \Pi - \nu \operatorname{curl}^2 \mathbf{u},
\end{equation}
where $\nu$ denotes the kinematic viscosity; the remaining equations~\eqref{eq:12-3} and \eqref{eq:12-4} are unaffected.

Consider the symmetric deformation of the cylinder specified by
\begin{equation}
\label{eq:12-60}
\varpi = R + \epsilon_0 e^{\sigma t + ikz},
\end{equation}
where
\begin{equation}
\label{eq:12-61}
\Pi = \epsilon_0 \Pi_0 I_0(k\varpi) e^{i kz + \sigma t},
\end{equation}
and
\begin{equation}
\label{eq:12-62}
\Pi = -\delta\mathfrak{B} + \delta p / \rho = -4\pi G\rho R_\epsilon K_0(x) I_0(k\varpi) e^{ikz+\sigma t} + \Pi_0 I_0(k\varpi) e^{ikz+\sigma t}.
\end{equation}
In the terminology of \S~61~(a), $\operatorname{grad}\Pi$ is a poloidal vector with the defining scalar
\begin{equation}
\label{eq:12-63}
\epsilon_0 \frac{\Pi_0}{k^2} I_0(k\varpi) e^{ikz+\sigma t}.
\end{equation}
From equations~\eqref{eq:12-59} and the poloidal character of $\operatorname{grad}\Pi$, we conclude that $\mathbf{u}$ is also poloidal; we can accordingly express it in the manner (cf.\ Chapter~VI~(201))
\begin{equation}
\label{eq:12-64}
\mathbf{u} = \left[-ik\varpi U + \frac{1}{k} \frac{d}{d\varpi}(\varpi U)\frac{1}{\varpi}\right] e^{ikz+\sigma t},
\end{equation}
where $U$ is a function of $\varpi$ only.

Remembering that the operation of $\operatorname{curl}$ on a poloidal axisymmetric vector is equivalent to the operation of $-\Delta_1$ on the defining scalar, we obtain from equation~\eqref{eq:12-59} the scalar equation
\begin{equation}
\label{eq:12-65}
\sigma U = -\epsilon_0 \frac{\Pi_0}{k} I_0(k\varpi) + \nu \left(\frac{d^2}{d\varpi^2} + \frac{1}{\varpi}\frac{d}{d\varpi} - k^2\right)U,
\end{equation}
or, rearranging, we have
\begin{equation}
\label{eq:12-66}
\frac{d^2 U}{d\varpi^2} + \frac{3}{\varpi}\frac{dU}{d\varpi} + \left(s_1^2 + \frac{1}{\varpi^2}\right)U = \epsilon_0 \frac{\Pi_0}{\nu k} \frac{I_0(k\varpi)}{\varpi^2}.
\end{equation}
Making use of the relation
\begin{equation}
\label{eq:12-67}
\left(\frac{d^2}{d\varpi^2} + \frac{3}{\varpi}\frac{d}{d\varpi}\right) \frac{I_1(\kappa\varpi)}{\varpi} = \kappa^2 \frac{I_1(\kappa\varpi)}{\varpi},
\end{equation}
we can readily write down the general solution of equation~\eqref{eq:12-66} which is free of singularity at the origin; it is given by
\begin{equation}
\label{eq:12-68}
U = \frac{1}{\varpi}\left[A \frac{I_1(\ell\varpi)}{\ell} + \frac{\epsilon_0 \Pi_0}{\nu k} \frac{I_1(k\varpi)}{k}\right] e^{ikz+\sigma t},
\end{equation}
where
\begin{equation}
\label{eq:12-69}
\ell^2 = k^2 + \sigma/\nu.
\end{equation}
and $A$ is a constant of integration. The corresponding solutions for the components of the velocity are
\begin{equation}
\label{eq:12-70}
u_\varpi = x \left[A I_0(\ell\varpi) - \frac{\epsilon_0 \Pi_0}{\nu k^2} I_0(k\varpi)\right] e^{ikz+\sigma t},
\end{equation}
and
\begin{equation}
\label{eq:12-71}
u_z = ik\left[A I_1(\ell\varpi) - \frac{\epsilon_0 \Pi_0}{\sigma} I_1(k\varpi)\right] e^{ikz+\sigma t}.
\end{equation}

The solution represented by equations~\eqref{eq:12-70} and \eqref{eq:12-71} must satisfy certain boundary conditions. These are: (1) the radial component of the velocity $u_\varpi$ must be compatible with the assumed form of the deformed boundary, namely,~\eqref{eq:12-60}; (2) the tangential viscous stresses must vanish at $\varpi = R$; and (3) the $(\varpi,\varpi)$-component of the total stress tensor must vanish on the deformed boundary.

The first of the foregoing conditions gives
\begin{equation}
\label{eq:12-72}
\epsilon_0 \sigma = x\left[A I_0(y) - \frac{\epsilon_0 \Pi_0}{\sigma} I_0(x)\right],
\end{equation}
where we have used the abbreviations
\begin{equation}
\label{eq:12-73}
x = kR \quad \text{and} \quad y = \ell R.
\end{equation}
The tangential viscous stress which must vanish at $\varpi = R$ is
\begin{equation}
\label{eq:12-74}
p_{\varpi z} = \rho\nu\left(\frac{\partial u_z}{\partial \varpi} + \frac{\partial u_\varpi}{\partial z}\right).
\end{equation}

With the representation~\eqref{eq:12-64} for $\mathbf{u}$,
\begin{equation}
\label{eq:12-75}
p_{\varpi z} = \rho\nu k \left[-ik\left(\frac{d}{d\varpi} + \frac{1}{\varpi}\right)U + ik^2 U\right] e^{ikz+\sigma t},
\end{equation}
and this must vanish at $\varpi = R$. For $U$ given by~\eqref{eq:12-68}, the condition gives
\begin{equation}
\label{eq:12-76}
A(k^2 + \tfrac{1}{2}\ell^2) I_0(y) - 2\frac{\epsilon_0 \Pi_0}{\sigma} k^2 I_0(x) = 0,
\end{equation}
or
\begin{equation}
\label{eq:12-77}
A = \frac{\epsilon_0 \Pi_0}{\sigma} \frac{2k^2}{k^2 + \frac{1}{2}\ell^2} \frac{I_0(x)}{I_0(y)}.
\end{equation}
Inserting this value of $A$ in equation~\eqref{eq:12-72}, we obtain
\begin{equation}
\label{eq:12-78}
\sigma^2 = \frac{\Pi_0}{\epsilon_0} \frac{k^2 - \ell^2}{k^2 + \frac{1}{2}\ell^2} x I_0(x).
\end{equation}

It remains to satisfy the last of the boundary conditions which requires the vanishing of the $(\varpi,\varpi)$-component of the total stress tensor on~\eqref{eq:12-60}. The required component of the stress tensor is
\begin{equation}
\label{eq:12-79}
P_{\varpi\varpi} = p_0 + \delta p - 2\rho\nu \frac{\partial u_\varpi}{\partial \varpi},
\end{equation}
where $p_0$ is the pressure in the unperturbed configuration given by~\eqref{eq:12-22}. Retaining only terms of the first order in $\epsilon$, we find that the vanishing of $P_{\varpi\varpi}$ on~\eqref{eq:12-60} requires
\begin{equation}
\label{eq:12-80}
-2\pi G\rho R_\epsilon e^{ikz+\sigma t} + \left(\frac{\delta p}{\rho}\right)_R - 2\nu\left(\frac{\partial u_\varpi}{\partial \varpi}\right)_R = 0.
\end{equation}
On the other hand, from equations~\eqref{eq:12-61} and \eqref{eq:12-62},
\begin{equation}
\label{eq:12-81}
-4\pi G\rho R_\epsilon K_0(x) I_0(x) e^{ikz+\sigma t} + \epsilon_0 \Pi_0 I_0(x) e^{ikz+\sigma t} = 0.
\end{equation}
Eliminating $(\delta p)_R$ between these last two equations, we obtain
\begin{equation}
\label{eq:12-82}
[4\pi G\rho R K_0(x) I_0(x) - \tfrac{1}{2}] + \Pi_0 I_0(x) \frac{d}{dx} = \frac{2\nu\rho}{\epsilon_0 k^2}.
\end{equation}
Now from equation~\eqref{eq:12-70}
\begin{equation}
\label{eq:12-83}
\left(\frac{\partial u_\varpi}{\partial \varpi}\right)_R = x \left[A\ell I'_0(y) - \frac{\epsilon_0 \Pi_0}{\sigma}k I'_0(x)\right] e^{ikz+\sigma t}.
\end{equation}
Using this result in equation~\eqref{eq:12-82}, we find
\begin{equation}
\label{eq:12-84}
\left(\frac{d\varpi}{\varpi}\right) = \epsilon_0 \Pi_0 I_0(x) + 4\pi G\rho R_\epsilon K_0(x) I_0(x) - \frac{\epsilon_0 \Pi_0}{2\nu k} I'_0(x).
\end{equation}
and substituting for $A$ in this equation its value~\eqref{eq:12-77}, we obtain
\begin{equation}
\label{eq:12-85}
\Pi_0 I_0(x) \left\{1 + \frac{2k^2\ell}{k^2+\frac{1}{2}\ell^2} \frac{I'_0(y)}{I_0(y)} \left[\frac{I_0(x)}{I'_0(x)} - \frac{I_0(y)}{I'_0(y)}\right]\right\} = -4\pi G\rho R_\epsilon [K_0(x)I_0(x) - \tfrac{1}{2}].
\end{equation}
Finally, eliminating $\Pi_0$ between equations~\eqref{eq:12-78} and \eqref{eq:12-85}, we obtain the characteristic equation
\begin{equation}
\label{eq:12-86}
\sigma^2 \left\{1 + \frac{2k^2\ell}{k^2+\frac{1}{2}\ell^2} \frac{I'_0(y)}{I_0(y)} \left[\frac{I_0(x)}{I'_0(x)} - \frac{I_0(y)}{I'_0(y)}\right]\right\} \frac{k^2+\frac{1}{2}\ell^2}{k^2-\ell^2} = -4\pi G\rho \frac{x I'_0(x)}{I_0(x)} [K_0(x)I_0(x)-\tfrac{1}{2}],
\end{equation}
where according to equations~\eqref{eq:12-69} and \eqref{eq:12-73}
\begin{equation}
\label{eq:12-87}
\sigma R^2/\nu = y^2 - x^2.
\end{equation}
Using this last relation, we can rewrite equation~\eqref{eq:12-86} in the form
\begin{equation}
\label{eq:12-88}
(x^2-y^2)\left\{1 - \frac{2xy}{x^2+\frac{1}{2}y^2} \frac{I_1(x)}{I_0(x)}\frac{I_0(y)}{I_1(y)} - \frac{1}{x^2+\frac{1}{2}y^2}\right\} = -\frac{J}{(y^2-x^2)^2} \frac{x I_1(x)}{I_0(x)} [K_0(x)I_0(x)-\tfrac{1}{2}],
\end{equation}
where
\begin{equation}
\label{eq:12-89}
J = 4\pi G\rho R^4/\nu^2.
\end{equation}
Equation~\eqref{eq:12-88} can be written more conveniently in the form
\begin{equation}
\label{eq:12-90}
2\sigma(x^2+y^2)\frac{I_1(x)}{I_0(x)}\left\{\frac{I_0(x)}{I_0(y)}\frac{I_1(y)}{I_1(x)} - \frac{2xy}{x^2+\frac{1}{2}y^2}\right\} [x^2-y^2] = -\frac{J}{x^2+y^2}[K_0(x)I_0(x)-\tfrac{1}{2}].
\end{equation}

One general result which follows at once from equation~\eqref{eq:12-90} is that \emph{the instability of the modes for $0 < x < x_a$} $(= 1.0668)$ \emph{is not affected by viscosity}.

\subsection*{(a) The case when the effects of viscosity are dominant}

Before we consider the characteristic equation~\eqref{eq:12-90} quite generally, it will be useful to examine the case when $\nu \to \infty$ and the effects of viscosity dominate inertia. In this limit, we may write
\begin{equation}
\label{eq:12-91}
y = x + \delta \quad \text{and treat } \delta \text{ as small.}
\end{equation}
Supposing this is the case, we find that to the first order in $\delta$, the quantity on the left-hand side of equation~\eqref{eq:12-90} becomes
\begin{equation}
\label{eq:12-92}
4x^3\delta - 4x^3\delta\frac{I_1^2}{I_1^2 - I_0 I_2},
\end{equation}
where the argument of all the Bessel functions is now $x$. By making use of the known properties of the Bessel functions, we can reduce~\eqref{eq:12-92} to the form
\begin{equation}
\label{eq:12-93}
4x^3\delta \left\{1 - \frac{x^2}{2}\left[\left(1 + \frac{1}{x^2}\right)I_1^2 - I_0 I_2\right]^{-1}\right\}.
\end{equation}
Since in the present approximation
\begin{equation}
\label{eq:12-94}
2\sigma\delta \approx \sigma R^2/\nu,
\end{equation}
the characteristic equation reduces in this limit to
\begin{equation}
\label{eq:12-95}
\sigma = -J \frac{x I_1}{I_0 \cdot x^2 [I_1^2/I_0^2-(1+x^2)^{-1}]} [K_0 I_0 - \tfrac{1}{2}],
\end{equation}
or substituting for $J$ its value~\eqref{eq:12-89}, we have
\begin{equation}
\label{eq:12-96}
\sigma = \frac{4\pi G\rho R^2}{\nu} \frac{K_0(x)I_0(x)-\frac{1}{2}}{x^2[I_1^2(x)/I_0^2(x)-(1+x^2)^{-1}]},
\end{equation}
where the arguments of the Bessel functions have been restored.

\begin{table}[ht]
\centering
\caption{The dispersion relation in the limit of high viscosity for the gravitationally unstable modes of an infinite cylinder}
\label{tab:LIV}
\begin{tabular}{cccccc}
\hline
$x$ & $\frac{4\pi G\rho R^2}{\nu}\sigma$ & $x$ & $\frac{4\pi G\rho R^2}{\nu}\sigma$ & $x$ & $\frac{4\pi G\rho R^2}{\nu}\sigma$ \\
\hline
0.01 & 0.1662 & 0.30 & 0.1517 & 0.70 & 0.0335 \\
0.05 & 0.1666 & 0.35 & 0.1462 & 0.75 & 0.0217 \\
0.10 & 0.1666 & 0.40 & 0.1379 & 0.80 & 0.0119 \\
0.15 & 0.1662 & 0.45 & 0.1269 & 0.85 & 0.0044 \\
0.20 & 0.1649 & 0.50 & 0.1135 & 0.90 & 0.0000 \\
0.25 & 0.1623 & 0.55 & 0.0979 & 0.95 & $-0.0019$ \\
& & 0.60 & 0.0805 & 1.00 & $-0.0019$ \\
& & 0.65 & 0.0568 & 1.0668 & 0 \\
\hline
\end{tabular}
\end{table}

\figurePlaceholder{124}{The dispersion relations in the limit of high viscosity for the modes of gravitational (curves labelled 1) and capillary (curves labelled 2) instability of an infinite cylinder. The abscissa measures the wave number in the unit $1/R$ and the ordinate measures the growth rate in the unit $4\pi G\rho R^2/\nu$ for the gravitational modes and in the unit $T/(\rho R^3)$ for the capillary modes.}

In the limit $x \to \infty$, the function of $x$ which occurs on the right-hand side of equation~\eqref{eq:12-96} given in the unit $4\pi G\rho R^2/\nu$. Values of this function for some values of $x$ are given in Table~LIV (see also Fig.~124). We observe that in contrast to the case when viscosity is absent, there is, in this limit, no finite mode of maximum instability: $\sigma$ is, in fact, a monotone decreasing function; and, moreover,
\begin{equation}
\label{eq:12-97}
\sigma \to -\frac{2\pi G\rho R^2}{\nu}(y + \tfrac{1}{2}\log x) \quad \text{as } x \to 0,
\end{equation}
where $\gamma$ $(= 0.5772\ldots)$ is Euler's constant. Thus, when viscosity is paramount, the cylinder will not have a tendency to break up at points separated by distances comparable to the circumference of the cylinder; rather, it will have a tendency to just give way at a few and distant places.

\begin{table}[ht]
\centering
\caption{Values of the function $\Phi(x)$}
\label{tab:LV}
\begin{tabular}{cccc}
\hline
$x$ & $\Phi(x)$ & $x$ & $\Phi(x)$ \\
\hline
0 & 0 & 0.5 & 0.07 \\
0.00002 & 0.00005 & 0.6 & 0.07 \\
0.0005 & 0.00050 & 0.7 & 0.06 \\
0.001 & 0.0031 & 0.8 & 0.05 \\
0.01 & 0.0076 & 0.9 & 0.03 \\
0.05 & 0.0406 & 0.95 & 0.02 \\
0.1 & 0.0549 & 1.0 & 0.01 \\
0.2 & 0.0605 & 1.05 & 0.002 \\
0.3 & 0.0645 & 1.0668 & 0 \\
0.4 & 0.0664 & & \\
\hline
\end{tabular}
\end{table}

The function
\begin{equation}
\label{eq:12-98}
\Phi(x) = \frac{x^2 I_1^2(x)/I_0^2(x) - (1+x^2)^{-1}}{x^2[I_1^2(x)/I_0^2(x)-(1+x^2)^{-1}]}
\end{equation}
which occurs in the denominator of the expression for $\sigma$ given by equation~\eqref{eq:12-96}, are given in Table~LV. It is seen that for $\sigma \leq x \leq x_a$, the function does not depart from its value at $x = 0$ (namely, 0) by more than one per cent. Therefore, to a good approximation, we can write
\begin{equation}
\label{eq:12-99}
\sigma = \frac{2\pi G\rho R^2}{\nu} [K_0(x)I_0(x)-\tfrac{1}{2}].
\end{equation}

\subsection*{(b) The general case}

We shall now show how one can explicitly obtain the dispersion relation implicitly expressed by equations~\eqref{eq:12-86} and \eqref{eq:12-87}. For an assigned value of $x$ $(< 1.0668)$, we can derive a $(\sigma, J)$-relation by evaluating $\sigma$ by equation~\eqref{eq:12-86} and $J$ by equation~\eqref{eq:12-89} for a sequence of values of $y$. Relations derived in this manner for a number of assigned values of $x$ are given in Table~LVI and exhibited in Fig.~125. By reading the curves in Fig.~125, or by interpolation among the tabulated values, we can deduce the values of $\sigma$ for different values of $x$ and a given value of $J$. Dispersion relations derived in this manner for different values of $J$ are illustrated in Fig.~126. The progressive shifting of the mode of maximum instability towards smaller values of $x$ as $J$ decreases is apparent.

\begin{table}[ht]
\centering
\caption{The $(\sigma, J)$-relations for various initially assigned values of $x$ ($\sigma$ is given in the unit $\nu/R^2$)}
\label{tab:LVI}
\begin{tabular}{cccccccccc}
\hline
\multicolumn{2}{c}{$x=0.7$} & \multicolumn{2}{c}{$x=0.2$} & \multicolumn{2}{c}{$x=0.5$} & \multicolumn{2}{c}{$x=0.4$} & \multicolumn{2}{c}{$x=0.5$} \\
$\sigma$ & $J$ & $\sigma$ & $J$ & $\sigma$ & $J$ & $\sigma$ & $J$ & $\sigma$ & $J$ \\
\hline
\multicolumn{10}{c}{(See original table for detailed values)} \\
\hline
\end{tabular}
\end{table}

\figurePlaceholder{125}{The effect of viscosity on the gravitationally unstable modes of an infinite cylinder. The $(\sigma,J)$-curves for different initially assigned values of the wave number $x$. The curves labelled 1, 2,\ldots, 9 are for values of $x = 0.9, 0.8, 0.7, 0.6, 0.5, 0.4, 0.3, 0.2, 0.1$, respectively.}

\figurePlaceholder{126}{The effect of viscosity on the gravitational instability of an infinite cylinder. The abscissa measures the wave number in the unit $1/R$ and the ordinate the growth rate in the unit $\nu/R^2$. The different curves are labelled by the values of $J$ $(= 4\pi G\rho R^4/\nu^2)$ to which they refer.}

\section{The effect of a uniform axial magnetic field on the gravitational instability of an infinite cylinder}\label{sec:110}
We shall now consider the effect of a uniform magnetic field along the axis of an infinite cylinder on its gravitational instability. We shall limit this present discussion to the case when the fluid is inviscid and of infinite electrical conductivity. The relevant perturbation equations are (cf.\ equations~IV~(10) and (31))
\begin{equation}
\label{eq:12-100}
\frac{\partial \mathbf{u}}{\partial t} = -\operatorname{grad} \Pi + \frac{\mu}{4\pi\rho}\frac{\partial \mathbf{h}}{\partial z},
\end{equation}
\begin{equation}
\label{eq:12-101}
\frac{\partial \mathbf{h}}{\partial t} = \operatorname{curl}(\mathbf{u} \times \mathbf{H}),
\end{equation}
\begin{equation}
\label{eq:12-102}
\operatorname{div} \mathbf{u} = 0, \quad \operatorname{div} \mathbf{h} = 0,
\end{equation}
and
\begin{equation}
\label{eq:12-103}
\Pi = -\delta\mathfrak{B} + \frac{\delta p}{\rho} + \frac{\mu H H_z}{4\pi\rho},
\end{equation}
where the various symbols have their usual meanings.

Since in the absence of a magnetic field the cylinder is unstable only for certain axisymmetric modes, we shall restrict our consideration of equations~\eqref{eq:12-100}--\eqref{eq:12-103} to the axisymmetric case.

As in \S~109, let the deformation of the surface be described by
\begin{equation}
\label{eq:12-104}
\varpi = R + \epsilon_0 e^{ikz+\sigma t}.
\end{equation}
The corresponding solution for $\Pi$ is (cf.\ equations~\eqref{eq:12-61} and \eqref{eq:12-62})
\begin{equation}
\label{eq:12-105}
\Pi = \epsilon_0 \Pi_0 I_0(k\varpi) e^{ikz+\sigma t};
\end{equation}
and the defining scalar of $\operatorname{grad}\Pi$ is
\begin{equation}
\label{eq:12-106}
\epsilon_0 \frac{\Pi_0}{k^2} I_0(k\varpi) e^{ikz+\sigma t},
\end{equation}
where $\Pi_0$ is a constant; and the defining scalar of $\operatorname{grad}\Pi$ is
\begin{equation}
\label{eq:12-107}
\epsilon_0 \frac{\Pi_0}{k^2} I_0(k\varpi) e^{ikz+\sigma t}.
\end{equation}
We shall express $\mathbf{u}$ in terms of a scalar function $U(\varpi)$ and $P(\varpi)$ in the manner (cf.\ equation~\eqref{eq:12-64}). Then $\mathbf{u} \times \mathbf{H}$ is a toroidal vector with the defining scalar
\begin{equation}
\label{eq:12-108}
ikHU(\varpi)e^{ikz+\sigma t},
\end{equation}
and $\operatorname{curl}(\mathbf{u} \times \mathbf{H})$ is a poloidal vector with the same defining scalar. From equation~\eqref{eq:12-101} it now follows that the defining scalar of $\mathbf{h}$ is
\begin{equation}
\label{eq:12-109}
\frac{ikH}{\sigma}U(\varpi)e^{ikz+\sigma t},
\end{equation}
and the components of $\mathbf{h}$ are given by
\begin{equation}
\label{eq:12-110}
h_\varpi^{(\mathrm{in})} = \frac{ik^2 H}{\sigma}\left(-ik U + \frac{1}{\varpi}\frac{d}{d\varpi}(\varpi U)\right)e^{ikz+\sigma t},
\end{equation}
and
\begin{equation}
\label{eq:12-111}
h_z^{(\mathrm{in})} = \frac{-k^2 H}{\sigma}U e^{ikz+\sigma t}.
\end{equation}

In writing equation~\eqref{eq:12-111}, we have distinguished $\mathbf{h}$ by $\mathbf{h}^{(\mathrm{in})}$ to emphasize that this is the perturbation in the magnetic field in the interior of the fluid. There will be a corresponding perturbation in the field exterior to the fluid; and since this latter perturbation is in a vacuum, it can be derived from a potential (in the same way as an irrotational velocity field) and can be expressed in the form
\begin{equation}
\label{eq:12-112}
\mathbf{h}^{(\mathrm{ext})} = A\{-K_1(k\varpi)\hat{e}_\varpi + iK_0(k\varpi)\hat{e}_z\}e^{ikz+\sigma t},
\end{equation}
where $A$ is a constant to be determined.

Returning to equation~\eqref{eq:12-100} and replacing each term in the equation by its defining scalar, we obtain
\begin{equation}
\label{eq:12-113}
\sigma U = -\epsilon_0 \frac{\Pi_0}{k} I_0(k\varpi) - \frac{\mu k^2 H^2}{4\pi\rho\sigma}U.
\end{equation}
Hence,
\begin{equation}
\label{eq:12-114}
U = -\frac{\epsilon_0 \Pi_0}{\sigma + \frac{\mu k^2 H^2}{4\pi\rho\sigma}} \frac{I_0(k\varpi)}{k},
\end{equation}
where
\begin{equation}
\label{eq:12-115}
\Omega_a = \sqrt{\frac{\mu H^2 k^2}{4\pi\rho}}
\end{equation}
is the Alfv\'en frequency. The solution represented by equations~\eqref{eq:12-111}, \eqref{eq:12-112}, and \eqref{eq:12-114} must satisfy certain boundary conditions. These are: (1) the radial component of the velocity at $\varpi = R$ must be compatible with the assumed form of the boundary; (2) the normal component of $\mathbf{h}$, at $\varpi = R$, must be continuous across the boundary; and (3) the normal component of the total stress tensor must be continuous across the deformed boundary.

The first of the foregoing conditions gives (cf.\ equations~\eqref{eq:12-64} and \eqref{eq:12-114})
\begin{equation}
\label{eq:12-116}
u_\varpi(R) = -\frac{\epsilon_0 \Pi_0 k}{\sigma^2+\Omega_a^2}I_0'(kR)e^{ikz+\sigma t} = \epsilon_0\sigma e^{ikz+\sigma t},
\end{equation}
where $x = kR$. Alternatively, we can write
\begin{equation}
\label{eq:12-117}
\sigma^2 + \Omega_a^2 = -\frac{\Pi_0}{\epsilon_0}xI_1(x).
\end{equation}
From equations~\eqref{eq:12-111}, \eqref{eq:12-113}, and \eqref{eq:12-114}, we find that the continuity of $h_\varpi$ at $\varpi = R$ gives
\begin{equation}
\label{eq:12-118}
h_\varpi^{(\mathrm{in})}(R) = -\frac{k^2 H}{\sigma^2+\Omega_a^2}\epsilon_0\Pi_0 I_1(x) = -AK_1(kR) = h_\varpi^{(\mathrm{ext})}(R).
\end{equation}
Hence,
\begin{equation}
\label{eq:12-119}
A = \frac{k^2 H}{\sigma^2+\Omega_a^2}\frac{\epsilon_0\Pi_0 I_1(x)}{K_1(x)}.
\end{equation}
Therefore (see equation~\eqref{eq:12-112}),
\begin{equation}
\label{eq:12-120}
h_z^{(\mathrm{ext})}(R) = \frac{ik^2 H}{\sigma^2+\Omega_a^2}\frac{\epsilon_0\Pi_0 I_1(x)}{K_1(x)}K_0(x)e^{ikz+\sigma t}.
\end{equation}

It remains to satisfy the last of the boundary conditions, namely, the continuity of the total normal stress on the deformed boundary~\eqref{eq:12-104}. On the other hand, from equation~\eqref{eq:12-105}, it follows that
\begin{equation}
\label{eq:12-121}
-2\pi G\rho R\epsilon + \left(\frac{\delta p}{\rho} + \frac{\mu H h_z^{(\mathrm{in})}}{4\pi\rho}\right)_R = \frac{1}{2}\Pi_0 I_0(x) e^{ikz+\sigma t}.
\end{equation}
where the first term on the left-hand side is the value of the unperturbed pressure on the boundary displaced by the amount $\epsilon_0 e^{ikz+\sigma t}$.

Eliminating $(\delta p)_R$ from equations~\eqref{eq:12-121} and \eqref{eq:12-122} and making use of equation~\eqref{eq:12-119}, we obtain
\begin{equation}
\label{eq:12-122}
-4\pi G\rho R K_0(x)I_0(x)\epsilon_0 + \Pi_0 I_0(x) + \frac{\mu H^2}{4\pi\rho}\frac{x^2 I_1(x)K_0(x)}{K_1(x)(\sigma^2+\Omega_a^2)}\epsilon_0\Pi_0 = 0.
\end{equation}
Eliminating $\Pi_0/\epsilon_0$ between equations~\eqref{eq:12-117} and \eqref{eq:12-122}, we have
\begin{equation}
\label{eq:12-123}
\frac{\Pi_0}{\epsilon_0}\left\{1 + \frac{x^2}{x^2+\frac{1}{2}y^2}\frac{I_0(x)}{K_1(x)} \cdot \frac{x I_1(x)K_0(x)}{I_0(x)}\right\} = -4\pi G\rho [K_0(x)I_0(x)-\tfrac{1}{2}].
\end{equation}
Equations~\eqref{eq:12-117} and \eqref{eq:12-123} now give
\begin{equation}
\label{eq:12-124}
\sigma^2 + \Omega_a^2 = \left\{1 + \frac{x^2}{(x^2+\frac{1}{2}y^2)} \frac{I_1(x)K_0(x)}{K_1(x)I_0(x)}\right\}^{-1} 4\pi G\rho \frac{x I_1(x)}{I_0(x)}[K_0(x)I_0(x)-\tfrac{1}{2}],
\end{equation}
or, making use of the relation
\begin{equation}
\label{eq:12-125}
\frac{I_1(x)}{I_0(x)} K_0(x) + \frac{I_0(x)}{K_1(x)} I_0(x) = \frac{1}{x},
\end{equation}
we finally obtain
\begin{equation}
\label{eq:12-126}
\sigma^2 = 4\pi G\rho \frac{xI_1(x)}{I_0(x)}\left[K_0(x)I_0(x)-\tfrac{1}{2}\right] - \frac{\mu H^2 k^2}{4\pi\rho} \frac{x^2 I_1(x)K_0(x)}{I_0(x)K_1(x)}.
\end{equation}
Defining
\begin{equation}
\label{eq:12-127}
H_0 = 4\pi\rho R_1\sqrt{G\rho},
\end{equation}
we can rewrite the characteristic equation~\eqref{eq:12-126} in the form
\begin{equation}
\label{eq:12-128}
\sigma^2 = 4\pi G\rho \frac{xI_1(x)}{I_0(x)}\left\{K_0(x)I_0(x)-\tfrac{1}{2} - \left(\frac{H}{H_0}\right)^2 \frac{x^2 K_0(x)}{I_0(x)K_1(x)}\right\}.
\end{equation}
The functions
\begin{equation}
\label{eq:12-129}
\frac{xI_1(x)}{I_0(x)}\left[K_0(x)I_0(x)-\tfrac{1}{2}\right] \quad \text{and} \quad \frac{x^3 I_1(x)K_0(x)}{I_0^2(x)K_1(x)}
\end{equation}
which occur in equation~\eqref{eq:12-128} are listed in Table~LVII. Also, it may be noted in this connexion that for $R = 250$ parsecs and $\rho = 2 \times 10^{-24}$~g/cm$^3$, $H_0 \approx 5 \times 10^{-6}$ gauss; this is the order of magnitude of the magnetic field believed to prevail in the spiral arms of galaxies.

\subsection*{(a) The nature of the stabilizing effect of the axial magnetic field}

From the sign of the term in $H^2$ in equation~\eqref{eq:12-128} it is apparent that the axial magnetic field will have the effect of increasing the wavelength at which instability occurs; but \emph{no magnetic field, however strong, can stabilize the cylinder for disturbances of all wavelengths: the gravitational instability of sufficiently long waves will persist.} The reason for this lies, of course, in the logarithmic singularity of $\Delta\mathfrak{B}$ for wavelengths tending to infinity; and we can never compensate for this by any increase in the magnetic energy which follows a deformation.

In general, the wavelength at which instability occurs will be determined by the root of the equation
\begin{equation}
\label{eq:12-130}
\frac{xI_1(x)}{I_0(x)}\left[K_0(x)I_0(x)-\tfrac{1}{2}\right] = \left(\frac{H}{H_0}\right)^2 \frac{x}{I_0(x)K_1(x)}.
\end{equation}
From the tabulation of the functions of $x$ which occur on the two sides of this equation in Table~LVII, it is clear that the equation allows, for

\begin{table}[ht]
\centering
\caption{The auxiliary functions}
\label{tab:LVII}
\begin{tabular}{cccccc}
\hline
$x$ & $\psi$ & $\chi$ & $x$ & $\psi$ & $\chi$ \\
\hline
\multicolumn{6}{c}{(See original table for values)} \\
\hline
\end{tabular}
\end{table}

an assigned value of $H/H_0$, a single positive root (exclusive of the root $x = 0$). If $x_a$ denotes this root, then
\begin{equation}
\label{eq:12-131}
\sigma^2 < 0 \quad \text{for } x > x_a \quad \text{and} \quad \sigma^2 > 0 \quad \text{for } 0 < x < x_a.
\end{equation}
The cylinder is, therefore, unstable for all varicose deformations with wavelengths exceeding $2\pi R/x_a$, where $x_a$ now depends on the strength of the impressed magnetic field through $H/H_0$. And, as in the absence of the field, there is a mode of maximum instability which occurs at a wave number, $x_m$, where the quantity on the right-hand side of equation~\eqref{eq:12-128} attains its maximum value.

In Table~LVIII, the values of $x_a$, $x_m$, and $\sigma/\sqrt{4\pi G\rho}$ at $x_m$, are listed for a few values of $H/H_0$.

\begin{table}[ht]
\centering
\caption{The wave numbers $x_a$, $x_m$ at which instability occurs and at which it is a maximum}
\label{tab:LVIII}
\begin{tabular}{cccc}
\hline
$H/H_0$ & $x_a$ & $x_m$ & $\sigma_{\max}/\sqrt{4\pi G\rho}$ \\
\hline
0 & 1.067 & 0.580 & 0.2455 \\
0.1 & 1.055 & 0.560 & 0.2450 \\
0.2 & 1.020 & 0.535 & 0.2420 \\
0.3 & 0.958 & 0.475 & 0.2360 \\
0.5 & 0.762 & 0.295 & 0.2070 \\
0.7 & 0.485 & 0.155 & 0.1410 \\
0.9 & 0.180 & 0.065 & 0.0524 \\
\hline
\end{tabular}
\end{table}

This table exhibits the strong tendency towards stabilization which the magnetic field exerts. This is shown, especially, by the very rapid decrease in $\sigma$ at its maximum, as the strength of the impressed magnetic field is increased. The correspondingly rapid increase in the characteristic times for break-up which this entails is really an exponential one, as we shall now show.

Since $x_a = 0.055$ already for $H = H_0$, for $H > H_0$ we can replace the Bessel functions which occur in equation~\eqref{eq:12-121} by the dominant terms in their respective series expansions for $x \to 0$. Thus, we obtain
\begin{equation}
\label{eq:12-132}
\tfrac{1}{2}\gamma + \log \tfrac{1}{2}x = -2\left(\frac{H}{H_0}\right)^2.
\end{equation}
Hence,
\begin{equation}
\label{eq:12-133}
x_a = 2\exp[-(y+\tfrac{1}{2}+2H^2/H_0^2)],
\end{equation}
or, numerically,
\begin{equation}
\label{eq:12-134}
x_a = 0.6811 e^{-2H^2/H_0^2} \quad (H > H_0).
\end{equation}

The corresponding solution for $x_m$ can be obtained by a similar consideration of equation~\eqref{eq:12-128}. For $x \ll 1$, this equation gives
\begin{equation}
\label{eq:12-135}
\frac{x\sigma^2}{4\pi G\rho} = -\tfrac{1}{2}x^2(y + \tfrac{1}{2} + \log \tfrac{1}{2}x) - x^4,
\end{equation}
The expression on the right-hand side attains its maximum value when
\begin{equation}
\label{eq:12-136}
y + \tfrac{1}{2}\log\tfrac{1}{2}x + \tfrac{1}{2} + 2H^2/H_0^2 = 0.
\end{equation}
Hence
\begin{equation}
\label{eq:12-137}
x_m = 2\exp[-(y+1) - 2H^2/H_0^2] = 0.4134 e^{-2H^2/H_0^2}.
\end{equation}
We find from equation~\eqref{eq:12-135} that the corresponding expression for $\sigma_{\max}$ is
\begin{equation}
\label{eq:12-138}
\sigma_{\max}^2 = \frac{(4\pi G\rho)}{4e^2} x_m^2 = 0.2055 e^{-4H^2/H_0^2}.
\end{equation}

Equations~\eqref{eq:12-134}, \eqref{eq:12-137}, and \eqref{eq:12-138} (valid for $H > H_0$) emphasize the facts which are already apparent from Table~LVIII.

In illustration of the foregoing results, we may consider the following numerical example. An infinite cylinder of radius 250 parsecs and density $2 \times 10^{-24}$~g/cm$^3$ will, in the absence of a magnetic field, break-up into pieces of length of the order of $2.7 \times 10^3$ parsecs in a time of the order of $1.0 \times 10^8$ years (see p.~520). If there should be an axial magnetic field

of strength $7.5 \times 10^{-6}$ gauss $(= 1.5H_0)$, the corresponding figures are $3.4 \times 10^5$ parsecs and $1.1 \times 10^{14}$ years. The fact that the characteristic time for break-up can be prolonged by such prodigious factors is of considerable astronomical interest.

\section{The capillary instability of a liquid jet}\label{sec:111}
We shall now consider the stability of a liquid jet held together in its equilibrium state by capillary forces. This is the classical problem of Plateau.

The analysis of the preceding sections can be adapted to this problem with only slight modifications. Thus, in the inviscid case, equations~\eqref{eq:12-2} and~\eqref{eq:12-3} governing departures from the equilibrium state are, of course, of general validity; only, in place of equation~\eqref{eq:12-4}, we now have
\begin{equation}
\label{eq:12-139}
\Pi = \delta p / \rho,
\end{equation}
since in the present case there is no additional contribution derived from a potential.

A normal mode of perturbation of the jet can again be specified by defining the shape of the boundary as by equation~\eqref{eq:12-5}. And it is clear, also, that the arguments leading to the solutions~\eqref{eq:12-15} and~\eqref{eq:12-18} and the result~\eqref{eq:12-21} are unaffected. However, the analysis following equation~\eqref{eq:12-21} which seeks to determine $\Pi_0$ is not valid now: for, in the unperturbed state the pressure $p_0$ inside the jet is a constant and is
\begin{equation}
\label{eq:12-140}
p_0 = T/R,
\end{equation}
where $T$ denotes the surface tension. (We are restricting our present discussion to the case when the jet is in free space and the external pressure can be ignored; in \S~(b) below we consider a case in which the situation is exactly the reverse.) According to equation~X~(24), the condition we must satisfy in the perturbed state is that on the deformed boundary,
\begin{equation}
\label{eq:12-141}
p_0 + \delta p = T\left(\frac{1}{R_1} + \frac{1}{R_2}\right),
\end{equation}
where $R_1$ and $R_2$ are the principal radii of curvature of the surface defined by equation~\eqref{eq:12-5}. Clearly,
\begin{equation}
\label{eq:12-142}
\frac{1}{R_1} + \frac{1}{R_2} = \frac{1}{R} - \frac{(1-m^2-x^2)}{R^2}\epsilon e^{i(kz+m\varphi)+\sigma t},
\end{equation}
\begin{equation}
\label{eq:12-143}
\left(\frac{\delta p}{\rho}\right)_R = \epsilon_0 \Pi_0 I_m(x) e^{i(kz+m\varphi)+\sigma t},
\end{equation}
where as according to equations~\eqref{eq:12-15} and \eqref{eq:12-139},
\begin{equation}
\label{eq:12-144}
\frac{\Pi_0}{k^2} = \frac{T}{\rho R^2 I_m(x)} \frac{x^2(x)}{I_m(x)} (1-m^2-x^2).
\end{equation}
Therefore,
\begin{equation}
\label{eq:12-145}
\Pi_0 I_m(x) = \frac{T}{\rho R^2}(1-m^2-x^2) \epsilon_0.
\end{equation}
Inserting this expression for $\Pi_0/\epsilon_0$ in equation~\eqref{eq:12-21},
\begin{equation}
\label{eq:12-146}
\sigma_m^2 = \frac{T}{\rho R^3}\frac{x I'_m(x)}{I_m(x)}(1-m^2-x^2);
\end{equation}
in particular,
\begin{equation}
\label{eq:12-147}
\sigma_0^2 = \frac{T}{\rho R^3}\frac{xI_1(x)}{I_0(x)}(1-x^2).
\end{equation}
These formulae were first derived by Rayleigh.

\begin{table}[ht]
\centering
\caption{The dispersion relation for the modes of capillary instability of an infinite cylinder}
\label{tab:LIX}
\begin{tabular}{cccc}
\hline
$x$ & $\sigma_0^2/[T/(\rho R^3)]$ & $\frac{xI_1(x)}{I_0(x)}\sqrt{1-x^2}$ & \\
\hline
\multicolumn{4}{c}{(See original table for values)} \\
\hline
\end{tabular}
\end{table}

From equations~\eqref{eq:12-146} and~\eqref{eq:12-147} it is apparent that
\[
\sigma_m^2 < 0 \quad \text{for all } m \neq 0 \quad \text{and} \quad x > 0
\]
\[
\sigma_0^2 < 0 \quad \text{for } x > 1 \quad \text{and} \quad \sigma_0^2 > 0 \quad \text{for } 0 < x < 1.
\]
In other words: \emph{the liquid jet is stable for all purely non-axisymmetric deformations; but it is unstable for symmetric varicose deformations with wavelengths exceeding the circumference of the cylinder.} This last result is due to Plateau.

The value of $\sigma_0^2/(T/\rho R^3)$ for $0 \leq x < 1$ are given in Table~LIX. It is seen that in this case the mode of maximum instability occurs for $x = 0.697$; the corresponding wavelength of the varicose deformation is $\lambda = 2\pi R/(0.697) \approx 9.01 \times 2R$. And the value of $\sigma_0$ at its maximum is $0.3433\sqrt{T/(\rho R^3)}$. According to this last result, a water jet of diameter 1~cm has a characteristic time of break-up of about one-eighth sec.

\subsection*{(a) The origin of the capillary instability}

As in \S~108~(a), we can trace the origin of the capillary instability of a liquid jet to the sign of the change in the potential energy consequent to a deformation.

For the deformation specified by equation~\eqref{eq:12-36}, the volume per unit length is $\pi R_1^2(1+\frac{1}{2}\delta_{\epsilon,m}\epsilon^2/R^2)$, where
\begin{equation}
\label{eq:12-148}
\delta_{\epsilon,m} \text{ has the meaning given in equation~\eqref{eq:12-47}}.
\end{equation}
From equation~\eqref{eq:12-148}, it is clear that to the second order in $\epsilon$ the radius $R_1$ of the unperturbed cylinder cannot be the same as the mean radius $R$ of the deformed cylinder: the two radii must, in fact, be connected by the relation
\begin{equation}
\label{eq:12-149}
R_1^2 = R^2 - \tfrac{1}{2}(1+\delta_{\epsilon,m})\epsilon^2.
\end{equation}
or
\begin{equation}
\label{eq:12-150}
R = R_1 + \tfrac{1}{4}(1+\delta_{\epsilon,m})\frac{\epsilon^2}{R_1}.
\end{equation}

Now the potential energy of a system arising from capillary forces is simply proportional to the total superficial area. For the deformation given by equation~\eqref{eq:12-36}, this area per unit length is:
\begin{equation}
\label{eq:12-151}
\text{superficial area per unit length} = 2\pi R + \tfrac{1}{2}(1+\delta_{\epsilon,m})(k^2 R^2 + m^2 - 1)\frac{\epsilon^2}{R},
\end{equation}
or, expressing $R$ in terms of $R_1$ we have:

superficial area per unit length $= 2\pi R_1 + \frac{1}{2}(1+\delta_{\epsilon,m})k^2\epsilon^2/R_1$,

where $x = kR$. The change in the surface energy $\Delta\mathfrak{S}$ per unit length consequent on the deformation is, therefore, given by
\begin{equation}
\label{eq:12-152}
\Delta\mathfrak{S} = \tfrac{1}{2}\pi T(1+\delta_{\epsilon,m})(k^2 R^2 + m^2 - 1)\frac{\epsilon^2}{R}.
\end{equation}

(It is not necessary to distinguish any longer between $R$ and $R_1$.)

Comparing equation~\eqref{eq:12-152} for $\Delta\mathfrak{S}$ with equation~\eqref{eq:12-146} for $\sigma_m^2$, we observe that the criterion for stability which follows from the conditions $\Delta\mathfrak{S} > 0$ and $\sigma_m^2 < 0$ is the same. Indeed, by combining~\eqref{eq:12-152} with the kinetic energy of the associated motions given by~\eqref{eq:12-55}, we can define a Lagrangian which will lead to the same formula for $\sigma_m^2$ as~\eqref{eq:12-146}.

\subsection*{(b) The capillary instability of a hollow jet}

In formulating the boundary condition in the manner~\eqref{eq:12-141}, we are neglecting the inertia of the fluid which may be outside the jet. This is clearly permissible in the case of a liquid jet in air; but there are cases in which this is not permissible and it is the inertia of the fluid outside, rather than that of the fluid inside, that is important. An example is provided by a jet of air injected into water. The stability of such a hollow jet can be discussed very similarly to solid jets. Indeed, it is clear that the corresponding formula for $\sigma_0^2$ will involve the $K$-functions in place of the $I$-functions. Thus, for the important axisymmetric modes, $m = 0$, we will have
\begin{equation}
\label{eq:12-153}
\sigma_0^2 = \frac{T}{R^3\rho}\frac{xK_1(x)}{K_0(x)}(1-x^2)
\end{equation}
instead of~\eqref{eq:12-147}. The square roots of the function of $x$ which occurs on the right-hand side of equation~\eqref{eq:12-153} is listed in Table~LX (see also Fig.~127). We observe that the mode of maximum instability occurs at a somewhat longer wavelength than for a solid jet of the same diameter; and the characteristic time for break-up is about $2\frac{1}{2}$ times shorter.

\begin{table}[ht]
\centering
\caption{The dispersion relation for the modes of capillary instability of a hollow jet}
\label{tab:LX}
\begin{tabular}{cccc}
\hline
$x$ & $\sqrt{\frac{xK_1}{K_0}(1-x^2)}$ & $x$ & $\sqrt{\frac{xK_1}{K_0}(1-x^2)}$ \\
\hline
\multicolumn{4}{c}{(See original table for values)} \\
\hline
\end{tabular}
\end{table}

\figurePlaceholder{127}{The dispersion relations for the modes of capillary instability ($x < 1$) of a solid (curves labelled 1) and a hollow (curves labelled 2) jet. The abscissa measures the wave number in the unit $1/R$ and the ordinate the growth rate in the unit $\sqrt{T/(\rho R^3)}$. The mode of maximum instability occurs at $x = 0.697$ for a solid jet and at $x = 0.484$ for a hollow jet.}

\subsection*{(c) The effect of viscosity}

The effect of viscosity on the capillary instability of a liquid jet can be discussed in the manner of \S~109: in fact, a large part of the analysis carries over without any modifications. Thus, except for the definition of $\Pi$ in equation~\eqref{eq:12-62}, the analysis up to and including equation~\eqref{eq:12-80} holds. The condition that $p_{\varpi\varpi}$ given by equation~\eqref{eq:12-79} should vanish on the deformed boundary is now replaced by (cf.\ equation~\eqref{eq:12-141})
\begin{equation}
\label{eq:12-154}
P_{\varpi\varpi}^{(\mathrm{def})} = (p+\delta p - 2\rho\nu \frac{\partial u_\varpi}{\partial \varpi})|_R = T\left(\frac{1}{R_1}+\frac{1}{R_2}\right).
\end{equation}

Making use of equation~\eqref{eq:12-142} (for the case $m = 0$) and equation~\eqref{eq:12-139}, we find that equation~\eqref{eq:12-82} is replaced by
\begin{equation}
\label{eq:12-155}
\left[\frac{T}{R^2\rho}(1-x^2) + \Pi_0 I_0(x)\right]\epsilon_0 e^{ikz+\sigma t} = 2\nu\left(\frac{\partial u_\varpi}{\partial \varpi}\right)_R.
\end{equation}

Continuing the analysis beyond this point exactly as in \S~109 (using equation~\eqref{eq:12-155} instead of equation~\eqref{eq:12-82})), we shall finally obtain (cf.\ equations~\eqref{eq:12-89} and \eqref{eq:12-90})
\begin{equation}
\label{eq:12-156}
2x(x^2+y^2) \frac{I_1(x)}{I_0(x)}\left\{\frac{I_0(x)}{I_0(y)}\frac{I_1(y)}{I_1(x)} - \frac{2xy}{x^2+\frac{1}{2}y^2}\right\}(x^2-y^2) = -\frac{J x I_1(x)}{I_0(x)}(1-x^2),
\end{equation}
where, now,
\begin{equation}
\label{eq:12-157}
J = \frac{T}{R\rho}\frac{TR}{\rho\nu^2}.
\end{equation}

\figurePlaceholder{128}{The effect of viscosity on the modes of capillary instability of an infinite cylinder. The $(\sigma,J)$-curves for different initially assigned values of the wave number $x$. The curves labelled 1, 2,\ldots, 9 are for values of $x = 0.9, 0.8, 0.7, 0.6, 0.5, 0.4, 0.3, 0.2, 0.1$, respectively.}

\figurePlaceholder{129}{The effect of viscosity on capillary instability of an infinite cylinder. The abscissa measures the wave number in the unit $1/R$ and the ordinate the growth rate in the unit $\nu/R^2$. The different curves are labelled by the values of $J$ $(= TR/\rho\nu^2)$ to which they refer.}

We observe that equation~\eqref{eq:12-156} differs from equation~\eqref{eq:12-90} only by a redefinition of $J$ and the replacement of $(K_0 I_0 - \frac{1}{2})$ (which is proportional to $\Delta\mathfrak{B}$) by $(1-x^2)$ (which is proportional to $\Delta\mathfrak{S}$). This identity between the two equations, apart from the differences mentioned, is the present analogue (in the context of a cylindrical geometry) of the theorem of \S~99.

When viscosity is paramount, the equation for $\sigma$ which is analogous to~\eqref{eq:12-96} is
\begin{equation}
\label{eq:12-158}
\sigma = \frac{T}{6\rho\nu R}\frac{x^2(1-x^2)}{x^2[I_1^2(x)/I_0^2(x)-(1+x^2)^{-1}]}.
\end{equation}
The dependence of $\sigma$ on $x$ predicted by this equation is exhibited in Table~LXI. We again observe that in this limit there is no finite mode of maximum instability (see Fig.~124); and this means that under the circumstances, the jet will not break-up into `drops' at mutual distances comparable to the circumference of the cylinder; rather, it will have a tendency to just give way at a few and distant places. A separation of a liquid jet into numerous drops, or varicosity suggesting such a separation, may be taken as an indication that `fluidity has been sufficient to bring inertia into play'.

\begin{table}[ht]
\centering
\caption{The dispersion relation in the limit of high viscosity for the modes of capillary instability of an infinite cylinder}
\label{tab:LXI}
\begin{tabular}{cccccc}
\hline
$x$ & $\sigma\rho\nu R/T$ & $x$ & $\sigma\rho\nu R/T$ & $x$ & $\sigma\rho\nu R/T$ \\
\hline
\multicolumn{6}{c}{(See original table for values)} \\
\hline
\end{tabular}
\end{table}

In view of the behaviour (exhibited in Table~LV) of the function of $x$ given by equation~\eqref{eq:12-158}, we may, to a sufficient approximation, write
\begin{equation}
\label{eq:12-159}
\sigma = \frac{T}{6\rho\nu R}(1-x^2).
\end{equation}

More generally, the dispersion relation implicitly expressed by equation~\eqref{eq:12-156} can be exhibited in the manner described in \S~109 in the case of gravitational instability. The $(\sigma, J)$-relations which follow from equations~\eqref{eq:12-87} and \eqref{eq:12-156} for different assigned values of $x$ are given in Table~LXII and illustrated in Fig.~128; and the dispersion relations for different $J$'s are shown in Fig.~129.

\begin{table}[ht]
\centering
\caption{The $(\sigma, J)$-relations for various initially assigned values of $x$ ($\sigma$ is given in the unit $\nu/R^2$)}
\label{tab:LXII}
\begin{tabular}{cccccccccc}
\hline
\multicolumn{10}{c}{(See original table for detailed values)} \\
\hline
\end{tabular}
\end{table}

\section{The effect of a uniform axial magnetic field on the capillary instability of a liquid jet}\label{sec:112}
We shall now consider the effect of a uniform magnetic field along the axis of a liquid jet on its capillary instability. If we limit ourselves to the case when the liquid is of zero viscosity and resistivity, the effect of the magnetic field on capillary instability can be deduced from the analysis of \S~110. Thus, apart from the definition of $\Pi$ in equations~\eqref{eq:12-104} and \eqref{eq:12-105}, the analysis up to and including equation~\eqref{eq:12-120} holds. But in place of equations~\eqref{eq:12-121} and \eqref{eq:12-122} we now have
\begin{equation}
\label{eq:12-160}
\left(\frac{\delta p}{\rho} + \frac{\mu H h_z^{(\mathrm{in})}}{4\pi\rho}\right)_R = \epsilon_0\frac{T}{\rho R^2}(1-x^2)e^{ikz+\sigma t};
\end{equation}
and
\begin{equation}
\label{eq:12-161}
\left(\frac{\delta p}{\rho} + \frac{\mu H h_z^{(\mathrm{ext})}}{4\pi\rho}\right)_R = \epsilon_0 \Pi_0 I_0(x) e^{ikz+\sigma t};
\end{equation}
instead of equation~\eqref{eq:12-124}. Continuing the analysis beyond this point exactly as in \S~110, we shall finally obtain (in place of equations~\eqref{eq:12-126} and \eqref{eq:12-128})
\begin{equation}
\label{eq:12-162}
\sigma^2 = \frac{T}{\rho R^3}\frac{xI_1(x)}{I_0(x)}(1-x^2) - \frac{\mu H^2}{4\pi\rho}\frac{x^2 I_1(x)K_0(x)}{I_0(x)K_1(x)}.
\end{equation}

The magnetic field has clearly the effect of increasing the wavelength at which capillary instability occurs; but unlike in the case of gravitational instability, the capillary instability can be suppressed completely by a magnetic field of sufficient strength. This difference in the two cases arises from the different behaviours of $\Delta\mathfrak{S}$ (which is proportional to $1-x^2$ and $\Delta\mathfrak{B}$ (which is proportional to $I_0 K_0 - \frac{1}{2}$) in the limit of long wavelengths. Specifically, since
\begin{equation}
\label{eq:12-163}
I_1(x) K_0(x) < \tfrac{1}{2} \quad \text{for } x \geq 0,
\end{equation}
\emph{magnetic fields of strength}
\begin{equation}
\label{eq:12-164}
\gamma > H_c = \frac{T}{4\pi\rho R^2}\frac{I_0(x_c)K_1(x_c)}{K_0(x_c)I_1(x_c)},
\end{equation}
\emph{will stabilize the jet for varicose deformations of all wavelengths.}

With the known surface tension of mercury ($T = 470$~dyne/cm), the minimum field according to~\eqref{eq:12-164} is 17~G: a very small field, indeed. However, we shall see below that the finite electrical conductivity of mercury alters this expectation completely.

\subsection*{(a) The effect of finite electrical conductivity}

Since the resistivity $\eta$ of metallic liquids such as mercury is very high ($\sim 10^4$ under normal terrestrial conditions), it is important to ascertain how the finiteness of the resistivity affects the possibility of suppressing the capillary instability of a liquid jet by an axial magnetic field.

When we allow for finite resistivity, the relevant perturbation equations are:
\begin{equation}
\label{eq:12-165}
\rho\frac{\partial \mathbf{u}}{\partial t} = -\operatorname{grad}\Pi + \frac{\mu H}{4\pi\rho}\frac{\partial \mathbf{h}}{\partial z},
\end{equation}
\begin{equation}
\label{eq:12-166}
\frac{\partial \mathbf{h}}{\partial t} = \operatorname{curl}(\mathbf{u}\times\mathbf{H}) + \eta\nabla^2\mathbf{h},
\end{equation}
\begin{equation}
\label{eq:12-167}
\Pi = \frac{\delta p}{\rho} + \frac{\mu H h_z}{4\pi\rho},
\end{equation}
and $\mathbf{u}$ and $\mathbf{h}$ are solenoidal vectors. In seeking solutions of the foregoing equations, we shall again restrict ourselves to axisymmetric perturbations and consider the normal mode of deformation specified by equation~\eqref{eq:12-104}.

The solution for $\Pi$ is again given by equation~\eqref{eq:12-105}; and in the representation of axisymmetric solenoidal vectors used in the preceding sections, the defining scalar of $\operatorname{grad}\Pi$ is (again, as before)
\begin{equation}
\label{eq:12-168}
\epsilon_0 \frac{\Pi_0}{k^2}I_0(k\varpi)e^{ikz+\sigma t}.
\end{equation}

We shall now represent $\mathbf{u}$ and $\mathbf{h}$ in terms of two functions $U(\varpi)$ and $P(\varpi)$ in the manner (cf.\ equation~\eqref{eq:12-64})
\begin{equation}
\label{eq:12-169}
\mathbf{u} = [-ik\varpi U_1, \frac{1}{\varpi}\frac{d}{d\varpi}(\varpi U_1)]
\end{equation}
and
\begin{equation}
\label{eq:12-170}
\mathbf{h}^{(\mathrm{in})} = [-ik\varpi \mathcal{U}_1, \frac{d}{d\varpi}(\varpi^2 T_1)\frac{1}{\varpi}],
\end{equation}
where we have distinguished $\mathbf{h}$ by a superscript `(in)' to indicate that this is the perturbation in the field in the interior of the fluid. (We must also consider the perturbation in the field exterior to the fluid; and we shall do this presently.)

Returning to equations~\eqref{eq:12-165} and \eqref{eq:12-166} and replacing each of the terms in the equations by its respective defining scalar, we obtain
\begin{equation}
\label{eq:12-171}
\sigma U = -\epsilon_0 \frac{\Pi_0}{k} I_0(k\varpi) + \frac{ik\mu H}{4\pi\rho}P,
\end{equation}
and
\begin{equation}
\label{eq:12-172}
\sigma P = ikHU + \eta\Delta_1 P,
\end{equation}
where, for functions with the chosen dependence on $z$,
\begin{equation}
\label{eq:12-173}
\Delta_1 = \frac{d^2}{d\varpi^2} + \frac{3}{\varpi}\frac{d}{d\varpi} - k^2.
\end{equation}
Eliminating $U$ between equations~\eqref{eq:12-171} and \eqref{eq:12-172}, we obtain
\begin{equation}
\label{eq:12-174}
\eta\Delta_1 P - \sigma P = ik\epsilon_0\frac{ikH}{\sigma}I_0(k\varpi) + ik\mu H\frac{P}{\sigma}.
\end{equation}
Inserting for $\Delta_1$ its expression~\eqref{eq:12-173} and rearranging, we obtain
\begin{equation}
\label{eq:12-175}
\left[\frac{d^2}{d\varpi^2} + \frac{3}{\varpi}\frac{d}{d\varpi} - (\sigma^2 + \Omega_a^2)/(\sigma\eta) - k^2\right]P = -\epsilon_0 \frac{kH}{\sigma\eta}\frac{I_0(k\varpi)}{\varpi},
\end{equation}
where $\Omega_a$ denotes the Alfv\'en frequency defined in equation~\eqref{eq:12-115}.

Making use of the relation~\eqref{eq:12-67}, we can at once write down the solution of equation~\eqref{eq:12-175} which is free of singularity at the origin; it is
\begin{equation}
\label{eq:12-176}
P = \epsilon_0 \frac{H}{} \left[A \frac{I_1(\ell\varpi)}{\ell\varpi} + \frac{\epsilon_0 kH}{\sigma} \frac{I_1(k\varpi)}{k\varpi}\right],
\end{equation}
where $A$ is a constant and
\begin{equation}
\label{eq:12-177}
s^2 = k^2 + \frac{1}{\eta}\left(\sigma + \frac{\Omega_a^2}{\sigma}\right).
\end{equation}

From equation~\eqref{eq:12-171}, we find that the corresponding solution for $U$ is
\begin{equation}
\label{eq:12-178}
U = -\frac{\epsilon_0\Pi_0}{\sigma k}I_0(k\varpi) + \frac{ik\mu H}{4\pi\rho\sigma}P.
\end{equation}

With $P$ given by equation~\eqref{eq:12-176}, the components of $\mathbf{h}^{(\mathrm{in})}$ are:
\begin{equation}
\label{eq:12-179}
h_\varpi^{(\mathrm{in})} = ik\left[A I_0(\ell\varpi) + \frac{\epsilon_0 kH}{\sigma}I_0(k\varpi)\right]e^{ikz+\sigma t},
\end{equation}
\begin{equation}
\label{eq:12-180}
h_z^{(\mathrm{in})} = -\left[A\ell I_1(\ell\varpi) + \frac{\epsilon_0 kH}{\sigma}k I_1(k\varpi)\right]e^{ikz+\sigma t}.
\end{equation}
The corresponding perturbation in the magnetic field exterior to the fluid is derivable from a potential (being a vacuum field) and we can write (cf.\ equation~\eqref{eq:12-112})
\begin{equation}
\label{eq:12-181}
h_\varpi^{(\mathrm{ext})} = ik_0 BK_0(k\varpi)e^{ikz+\sigma t}
\end{equation}
and
\begin{equation}
\label{eq:12-182}
h_z^{(\mathrm{ext})} = ik_0 BK_1(k\varpi)e^{ikz+\sigma t},
\end{equation}
where $B$ is a constant.

Finally, we may note that the radial component of $\mathbf{u}$ is given by (cf.\ equations~\eqref{eq:12-169} and \eqref{eq:12-178}),
\begin{equation}
\label{eq:12-183}
u_\varpi = \left[-\frac{\epsilon_0\Pi_0}{\sigma}k I_1(k\varpi) - \frac{\Omega_a^2}{\sigma^2}[AI_0(\ell\varpi) + \cdots]\right]e^{ikz+\sigma t}.
\end{equation}

The solution represented by equations~\eqref{eq:12-179}--\eqref{eq:12-183} must satisfy certain boundary conditions. These are: (1) the velocity $u_\varpi$ given by equation~\eqref{eq:12-183} must be compatible, at $\varpi = R$, with the assumed form of the deformed boundary; (2) the field $\mathbf{h}$ must be continuous across the boundary; and (3) the normal component of the total stress must also be continuous across the boundary.

The first of the foregoing conditions gives
\begin{equation}
\label{eq:12-184}
\frac{\epsilon_0 \Pi_0}{\sigma^2+\Omega_a^2}xI_1(x) - A\frac{\Omega_a^2}{\sigma^2+\Omega_a^2}I_0(y) = -\epsilon_0\sigma,
\end{equation}
where
\begin{equation}
\label{eq:12-185}
x = kR \quad \text{and} \quad y = \ell R.
\end{equation}
And the continuity of $\mathbf{h}$ across the boundary requires the components of $\mathbf{h}$ given by equations~\eqref{eq:12-179} and \eqref{eq:12-180} to agree with those given by equations~\eqref{eq:12-181} and \eqref{eq:12-182} at $\varpi = R$. These conditions give
\begin{equation}
\label{eq:12-186}
B = -\frac{k}{K_1(x)}\left[AI_0(y) + \frac{\epsilon_0 kH}{\sigma}I_0(x)\right].
\end{equation}
From the equality of the two different expressions for $B$ in~\eqref{eq:12-186}, we find
\begin{equation}
\label{eq:12-187}
A(I_0(y)K_1(x) + I_1(y)K_0(x)) = -\frac{\epsilon_0 kH}{\sigma}(I_0(x)K_1(x)+I_1(x)K_0(x)).
\end{equation}
Thus,
\begin{equation}
\label{eq:12-188}
A = -\frac{\epsilon_0 kH}{\sigma}\frac{\Phi(x)}{\Phi(y)},
\end{equation}
where
\begin{equation}
\label{eq:12-189}
\Phi = y I_0(y)K_1(x) + K_0(x)I_1(y) = -1.
\end{equation}

With $A$ given by equation~\eqref{eq:12-188}, equation~\eqref{eq:12-184} becomes
\begin{equation}
\label{eq:12-190}
\frac{\Pi_0}{\epsilon_0} = -\frac{\sigma^2+\Omega_a^2}{xI_1(x)}\left\{1 + \frac{\Omega_a^2}{\sigma^2+\Omega_a^2}\frac{\Phi(x)}{\Phi(y)}I_0(y)\right\}.
\end{equation}

It remains to satisfy the last of the boundary conditions, namely, the continuity of the normal component of the total stress on the deformed surface. This condition leads to the same equations~\eqref{eq:12-160} and \eqref{eq:12-161}; from these equations we obtain (by subtraction)
\begin{equation}
\label{eq:12-191}
\epsilon_0\frac{T}{\rho R^2}(1-x^2) = \epsilon_0\Pi_0 I_0(x) + \frac{\mu H}{4\pi\rho}[h_z^{(\mathrm{ext})}(R) - h_z^{(\mathrm{in})}(R)].
\end{equation}

Now substituting for $h_z^{(\mathrm{ext})}$ $(= h_z^{(\mathrm{in})})$ from equation~\eqref{eq:12-180}, and rearranging, we obtain
\begin{equation}
\label{eq:12-192}
\frac{T}{\rho R^2}(1-x^2) = \Pi_0 I_0(x) + \frac{\mu H^2 k^2}{4\pi\rho\sigma}A\Omega_a \frac{I_1(y)}{y}.
\end{equation}

Inserting for $A$ in this last equation its value~\eqref{eq:12-188}, we obtain
\begin{equation}
\label{eq:12-193}
\frac{T}{\rho R^2}(1-x^2) = \Pi_0 I_0(x) - \frac{\Omega_a^2}{\sigma}\frac{\Phi(x)}{\Phi(y)}\frac{xI_1(y)}{y}.
\end{equation}

Finally, eliminating $\Pi_0/\epsilon_0 R$ between equations~\eqref{eq:12-190} and \eqref{eq:12-193}, we find
\begin{equation}
\label{eq:12-194}
\sigma^2\left[1 + \frac{\Omega_a^2}{\sigma^2}\frac{\Phi(x)}{\Phi(y)}\right](1-x^2) = \frac{T}{\rho R^3}\frac{xI_1(x)}{I_0(x)}(1-x^2),
\end{equation}
where it may be recalled that (see equation~\eqref{eq:12-177})
\begin{equation}
\label{eq:12-195}
y^2 = x^2 + \frac{R^2}{\eta}\left(\sigma + \frac{\Omega_a^2}{\sigma}\right).
\end{equation}

In the limit of zero resistivity,
\begin{equation}
\label{eq:12-196}
y \to \infty, \quad \Phi \to y I_0(y) K_1(x) \quad (\eta \to 0)
\end{equation}
and equation~\eqref{eq:12-194} tends to
\begin{equation}
\label{eq:12-197}
\sigma^2 = \frac{T}{\rho R^3}\frac{xI_1(x)}{I_0(x)}(1-x^2) - \Omega_a^2 \frac{x I_1(x)K_0(x)}{I_0(x)K_1(x)} \quad (\eta \to 0),
\end{equation}
in agreement with equation~\eqref{eq:12-162}. On the other hand, in the limit of infinite resistivity,
\begin{equation}
\label{eq:12-198}
y \to x, \quad \Phi \to 1,
\end{equation}
and
\begin{equation}
\label{eq:12-199}
\sigma^2 = \frac{T}{\rho R^3}\frac{xI_1(x)}{I_0(x)}(1-x^2) \quad (\eta \to \infty);
\end{equation}
in this limit the magnetic field has no effect on the capillary instability as should, indeed, be the case.

\subsection*{(b) The case of high resistivity}

When the resistivity $\eta$ is high but not infinite, we may write
\begin{equation}
\label{eq:12-200}
y = x + \delta \quad \text{and treat } \delta \text{ as a small quantity.}
\end{equation}
To the first order in $\delta$, equation~\eqref{eq:12-194} becomes
\begin{equation}
\label{eq:12-201}
\sigma^2\left[1 + \delta \frac{x^2}{\Phi_1^2\Phi_2}\frac{I_1}{x}\log(I_0) - \frac{I_0}{I_0} \frac{I_1}{x} \log(I_0)\right](1-x^2),
\end{equation}
where all the Bessel functions have now the argument $x$. An equivalent form of equation~\eqref{eq:12-201} is
\begin{equation}
\label{eq:12-202}
\sigma^2 \left[1 + \delta\frac{x^2}{I_0^2}\left\{\frac{I_0}{x}\log(I_0) - \frac{I_1}{I_0}\log I_1\right\}\right] = \frac{T}{\rho R^3}\frac{x I_1(x)}{I_0(x)}(1-x^2).
\end{equation}
By making use of the recurrence relations satisfied by the Bessel functions, we find
\begin{equation}
\label{eq:12-203}
\frac{1}{x}\log(xI_1) - \frac{I_0+xI_1}{I_1} = \frac{1}{x} - \frac{I_0-I_2}{2I_1} = \frac{1}{2} - \frac{I_2}{2I_1}.
\end{equation}
Moreover, in the present approximation (cf.\ equations~\eqref{eq:12-195} and \eqref{eq:12-198})
\begin{equation}
\label{eq:12-204}
2\delta\sigma = \frac{R^2}{\eta}(\sigma^2 + \Omega_a^2).
\end{equation}

We thus obtain
\begin{equation}
\label{eq:12-205}
\sigma^2\left[1 + \frac{\mu H^2 R}{4\pi\rho\eta\sigma}\frac{xI_1(x)}{I_0(x)}\left\{\frac{1}{2} - \frac{I_2}{2I_1}\right\}\right] = \frac{T}{\rho R^3}\frac{xI_1(x)}{I_0(x)}(1-x^2).
\end{equation}

In considering equation~\eqref{eq:12-205}, it is convenient to measure $\sigma$ in the unit
\begin{equation}
\label{eq:12-206}
\sigma_0 = \sqrt{\frac{T}{\rho R^3}} \sec^{-1};
\end{equation}

\begin{table}[ht]
\centering
\caption{The function $f(x) = \frac{xI_1(x) - I_0 I_1}{I_0 I_1}$}
\label{tab:LXIII}
\begin{tabular}{cccc}
\hline
$x$ & $f(x)$ & $x$ & $f(x)$ \\
\hline
0 & 0 & 0.5 & 0.0060 \\
0.1 & 0.000004 & 0.6 & 0.0088 \\
0.2 & 0.000064 & 0.7 & 0.0128 \\
0.3 & 0.000479 & 0.8 & 0.0175 \\
0.4 & 0.00170 & 0.9 & 0.0231 \\
& & 0.95 & 0.0272 \\
& & 1.0 & 0.0315 \\
\hline
\end{tabular}
\end{table}

equation~\eqref{eq:12-205} then takes the form
\begin{equation}
\label{eq:12-207}
\sigma^2 + 2\frac{\mu H^2 R^3}{4\pi\rho\eta}\sigma = \frac{xI_1(x)}{I_0(x)}(1-x^2).
\end{equation}
where
\begin{equation}
\label{eq:12-208}
\mathfrak{z} = \frac{\mu H^2 R^4}{4\pi\eta T \rho R^4} \frac{xI_1-I_0 I_1}{I_0 I_1}.
\end{equation}
For the validity of the present approximation, it is clearly necessary that the second term in parentheses on the left-hand side of equation~\eqref{eq:12-207} is small compared to unity. Assuming that this is the case, we obtain from equation~\eqref{eq:12-207} the following approximate formula for the roots
\begin{equation}
\label{eq:12-209}
\sigma = \pm\sqrt{f(x)} \cdot \sigma_0 - \frac{\mathfrak{z}}{2}\sigma_0.
\end{equation}

The function of $x$ which occurs as the factor of $\mathfrak{z}$ in this equation is listed briefly in Table~LXIII. This function, multiplied by $\mathfrak{z}$, must be subtracted from the corresponding entries in Table~LIX. From the relative values of the two functions listed in Tables~LIX and LXIII, we find that a value of $\mathfrak{z} \sim 6$ is needed to shift the critical wave number at which capillary instability occurs from 1.0 to 0.8. From this it would appear that values of $\mathfrak{z}$ as high as 20 or more will be needed to overcome significantly the paramount effects of finite resistivity.

With the values of the physical constants appropriate for mercury and liquid sodium, we find
\begin{equation}
\label{eq:12-210}
\mathfrak{z} = 1.33 \times 10^{-7}H^2R^4 \quad \text{for Hg with } T = 470 \text{ dyne/cm,}
\end{equation}
\[
\eta = 7.5 \times 10^3 \text{ cm}^2/\text{sec}, \quad \text{and } \rho = 13.5 \text{ g/cm}^3;
\]
\begin{equation}
\label{eq:12-211}
\text{and} \quad \mathfrak{z} = 7.4 \times 10^{-4}H^2R^4 \quad \text{for Na with } T = 200 \text{ dyne/cm,}
\end{equation}
\[
\eta = 7.6 \times 10^2 \text{ cm}^2/\text{sec}, \quad \text{and } \rho = 1 \text{ g/cm}^3.
\]
We might conclude from these values that for experiments with mercury, for example, magnetic field strengths of the order of $10^4$ gauss, or more, will be needed to demonstrate the stabilizing effect of an axial magnetic field.

\subsection*{(c) The general case}

While the first order formulae of \S~(b) would suggest that under ordinary laboratory conditions it will be difficult to overcome the effects of finite resistivity, it may still be useful to outline how one might treat equations~\eqref{eq:12-194} and \eqref{eq:12-195} generally.

Measuring $H$ and $\sigma$ in the units specified in equations~\eqref{eq:12-162} and \eqref{eq:12-206}, we first rewrite equations~\eqref{eq:12-194} and \eqref{eq:12-195} in the non-dimensional forms
\begin{equation}
\label{eq:12-212}
\sigma^2\left[1 + \frac{H^2 xI_1(y)}{I_0(y)\sigma\Phi_0 I_0(x)}\right](1-x^2) = \frac{xI_1(x)}{I_0(x)}\left[1 + \frac{H^2 x I_1(y)}{I_0(y)\sigma^2 \Phi_0 I_0(x)}\right](1-x^2),
\end{equation}
and
\begin{equation}
\label{eq:12-213}
y^2 = x^2 + \frac{S}{} (\sigma^2 + x^2 H^2 S),
\end{equation}
where
\begin{equation}
\label{eq:12-214}
S = \frac{R^2}{\eta}\sqrt{\frac{T}{\rho R^3}} = \frac{1}{\sqrt{\nu R}}\sqrt{\frac{T}{\rho}}.
\end{equation}

For a given set of values of $H$, $x$, and $y$, equation~\eqref{eq:12-212} can be solved as a quadratic for $\sigma^2$; and equation~\eqref{eq:12-213} can then be used to determine $S$. Keeping $H$ and $x$ fixed and allowing to run through a sequence of values, we can then derive a $(\sigma, S)$-relation (to be associated with the chosen values of $H$ and $x$). Interpolating among the $(\sigma, S)$-relations for different values of $x$, but the same value of $H$, we can obtain the $(\sigma, x)$-relation appropriate to some fixed value of $S$ and the chosen value of $H$. In this manner, the dispersion relation for any set of initial conditions can be deduced.

\section{The stability of the simplest solution of the equations of hydromagnetics}\label{sec:113}
Some interesting examples in the stability of cylindrical systems arise in connexion with the simplest solution (cf.\ Chapter~IV, \S~40~(a)),
\begin{equation}
\label{eq:12-215}
\mathbf{u}_s = \pm H_s\sqrt{\frac{\mu}{4\pi\rho}} \quad \text{and} \quad \Pi = \frac{p}{\rho} + \mathfrak{B} + \frac{\mu|\mathbf{H}|^2}{8\pi\rho} = \text{constant},
\end{equation}
of the equations governing an incompressible inviscid fluid of infinite electrical conductivity. Before we consider these examples, we shall first establish the stability of the solution~\eqref{eq:12-215} generally.

Writing
\begin{equation}
\label{eq:12-216}
h_s = \pm\sqrt{\frac{\mu}{4\pi\rho}}\Pi_s,
\end{equation}
we have the equations
\begin{equation}
\label{eq:12-217}
\frac{\partial u_j}{\partial t} + \frac{\partial}{\partial x_j}(u_s u_j - h_s h_j) = -\frac{\partial \Pi}{\partial x_j},
\end{equation}
\begin{equation}
\label{eq:12-218}
\frac{\partial h_j}{\partial t} + \frac{\partial}{\partial x_s}(h_s u_j - u_s h_j) = 0.
\end{equation}
In addition we have the equations expressing the solenoidal characters of $\mathbf{u}$ and $\mathbf{h}$.

Let
\begin{equation}
\label{eq:12-219}
u_s = h_s = U_s, \quad \frac{\partial U_s}{\partial x_s} = 0, \quad \text{and} \quad \Pi = \text{constant}
\end{equation}
represent the initial stationary state; and let
\begin{equation}
\label{eq:12-220}
U_s + \xi_s, \quad U_s + \zeta_s, \quad \text{and} \quad \Pi + P
\end{equation}
represent a slightly perturbed state. The equations governing the perturbations $\xi_s$, $\zeta_s$, and $P$ are
\begin{equation}
\label{eq:12-221}
\frac{\partial \xi_s}{\partial t} - \frac{\partial}{\partial x_j}(U_s\xi_j + U_j\xi_s) = -\frac{\partial P}{\partial x_s},
\end{equation}
and
\begin{equation}
\label{eq:12-222}
\frac{\partial \zeta_s}{\partial t} + \frac{\partial}{\partial x_j}(U_s\zeta_j - U_j\zeta_s) = 0,
\end{equation}
where
\begin{equation}
\label{eq:12-223}
X_j = h_j - u_j \quad \text{and} \quad Y_j = h_j + u_j.
\end{equation}
Let equations~\eqref{eq:12-221} and \eqref{eq:12-222} allow a stationary solution of the form
\begin{equation}
\label{eq:12-224}
X_j = \alpha_0 U_j \quad \text{and} \quad Y_j = \alpha_1 U_j,
\end{equation}
where $\alpha_0$ and $\alpha_1$ are constants.

From equations~\eqref{eq:12-221} and \eqref{eq:12-222} we obtain, by addition and subtraction the symmetrical pair of equations
\begin{equation}
\label{eq:12-225}
\frac{\partial Y_j}{\partial t} + X_s\frac{\partial Y_j}{\partial x_s} + Y_s\frac{\partial X_j}{\partial x_s} = -\frac{\partial\Pi}{\partial x_j},
\end{equation}
and
\begin{equation}
\label{eq:12-226}
\frac{\partial X_j}{\partial t} + Y_s\frac{\partial X_j}{\partial x_s} + X_s\frac{\partial Y_j}{\partial x_s} = -\frac{\partial\Pi}{\partial x_j}.
\end{equation}

In addition,
\begin{equation}
\label{eq:12-227}
\frac{\partial \xi_j}{\partial x_j} = 0.
\end{equation}

For solutions with a dependence on time given by
\begin{equation}
\label{eq:12-228}
e^{\sigma t},
\end{equation}
equation~\eqref{eq:12-226} becomes
\begin{equation}
\label{eq:12-229}
\sigma \zeta_j = -2U_s\frac{\partial \xi_j}{\partial x_s} + \frac{\partial P}{\partial x_j}.
\end{equation}

By multiplying equation~\eqref{eq:12-229} by $\xi_j^*$ and summing over $j$ as the notation implies and integrating over the volume $V$ occupied by the fluid, we obtain
\begin{equation}
\label{eq:12-230}
\sigma \int |\xi_j|^2 \, dV = -2 \int \xi_j^* U_s \frac{\partial \xi_j}{\partial x_s} \, dV + \int \frac{\partial}{\partial x_j}(P\xi_j^*) \, dV,
\end{equation}
\begin{equation}
\label{eq:12-231}
= -2 \int \xi_j^* U_s\frac{\partial \xi_j}{\partial x_s}\,dV + \int_S P\xi_j^* \, dS_j,
\end{equation}
where $S$ is the surface bounding $V$. We shall suppose that the conditions on the boundary are such that the surface integral in~\eqref{eq:12-230} vanishes. (It will suffice if either $P$ or the normal component of $\boldsymbol\xi$ vanishes on $S$.) Equation~\eqref{eq:12-230} will then give
\begin{equation}
\label{eq:12-232}
\sigma \int |\xi_j|^2 \, dV = -2 \int \xi_j^* U_s\frac{\partial \xi_j}{\partial x_s}\,dV.
\end{equation}
The complex conjugate of this equation is
\begin{equation}
\label{eq:12-233}
-\sigma^* \int |\xi_j|^2 \, dV = -2 \int \xi_j U_s \frac{\partial \xi_j^*}{\partial x_s}\,dV.
\end{equation}
Adding equations~\eqref{eq:12-232} and \eqref{eq:12-233}, we obtain
\begin{equation}
\label{eq:12-234}
(\sigma - \sigma^*) \int |\xi_j|^2 \, dV = -2\int (U_s\xi_j^*\frac{\partial \xi_j}{\partial x_s} + U_s\xi_j\frac{\partial \xi_j^*}{\partial x_s})\,dV = -\frac{\partial}{\partial x_s}\int |\xi_j|^2 U_s\, dV.
\end{equation}
Again, we shall suppose that the initial conditions and/or the boundary conditions are such that the surface integral in~\eqref{eq:12-234} vanishes. (It will suffice if the normal component of $\mathbf{U}$ vanishes on $S$.) Equation~\eqref{eq:12-234} will then give
\begin{equation}
\label{eq:12-235}
\mathrm{im}(\sigma) \int |\xi_j|^2 \, dV = 0.
\end{equation}
It follows that $\sigma$ is real; and \emph{this establishes the stability of the solution~\eqref{eq:12-215}}.

With the knowledge that $\sigma$ is real, we can derive a variational basis for its determination. For this purpose, we write
\begin{equation}
\label{eq:12-236}
\xi_j = \xi_j^{(R)} + i\xi_j^{(I)} \quad \text{and} \quad P = P^{(R)} + i\sigma P^{(I)}
\end{equation}
and separate the real and the imaginary parts of equation~\eqref{eq:12-229}. We then obtain the pair of equations
\begin{equation}
\label{eq:12-237}
\sigma\xi_j^{(R)} = -2U_s\frac{\partial \xi_j^{(I)}}{\partial x_s} + \frac{\partial P^{(R)}}{\partial x_j},
\end{equation}
and
\begin{equation}
\label{eq:12-238}
\sigma\xi_j^{(I)} = +2U_s\frac{\partial \xi_j^{(R)}}{\partial x_s} - \frac{\partial P^{(I)}}{\partial x_j}.
\end{equation}
(Both $\boldsymbol\xi^{(R)}$ and $\boldsymbol\xi^{(I)}$ are, of course, solenoidal.)

By multiplying equation~\eqref{eq:12-238} by $\xi_j^{(R)}$ and integrating over the volume of the fluid, we obtain
\begin{equation}
\label{eq:12-239}
\sigma^2 \int |\xi_j^{(R)}|^2 \, dV = 2 \int \xi_j^{(R)} U_s \frac{\partial \xi_j^{(R)}}{\partial x_s} \, dV - \int P^{(R)} \frac{\partial \xi_j^{(R)}}{\partial x_j}\, dV
\end{equation}
\begin{equation}
\label{eq:12-240}
= -2 \int \xi_j^{(R)} U_s \frac{\partial \xi_j^{(R)}}{\partial x_s}\,dV + \int \frac{\partial}{\partial x_j}(P^{(R)}\xi_j^{(R)}) U_s\,dS_j - \int P^{(R)}\frac{\partial}{\partial x_j}\xi_j^{(R)} \, dV
\end{equation}
\begin{equation}
\label{eq:12-241}
= -2 \int \xi_j^{(R)} U_s \frac{\partial \xi_j^{(R)}}{\partial x_s}\,dV + \int_S P^{(R)}\xi_j^{(R)} U_s\,dS_j - \int |\xi_j^{(R)}|^2 U_s^2 \, dV.
\end{equation}

If we suppose that the conditions on the boundary are such that the various surface integrals which have occurred during the reductions vanish, then
\begin{equation}
\label{eq:12-242}
\int |\xi_j^{(I)}|^2\,dV = 0.
\end{equation}

It can be shown in the usual manner that equation~\eqref{eq:12-242} provides the basis for a variational formulation of the problem. However, it is important to note that the extremal property of $\sigma$ is true only if, when performing the variation of the quantity on the right-hand side of equation~\eqref{eq:12-242}), we treat equation~\eqref{eq:12-237} as a constraint which must be satisfied by the variations as well.

\subsection*{(a) An example}

To illustrate in a concrete manner how the stability of the equipartition solutions arises, we shall consider an example for which $\sigma^2$ can be explicitly found. The stationary solution we consider, in cylindrical polar coordinates, is
\begin{equation}
\label{eq:12-243}
\mathbf{u} = \mathbf{H}\sqrt{\frac{\mu}{4\pi\rho}} = \mathbf{U} = U\left(0, \frac{\varpi}{p}, 1\right),
\end{equation}
where $U$ is a constant. We shall further suppose that the fluid is confined in a cylinder of radius $R$.

We have seen that in general the perturbation equations can be reduced to the form~\eqref{eq:12-229}. In cylindrical polar coordinates this equation becomes:
\begin{equation}
\label{eq:12-244}
\frac{\partial\xi_\varpi}{\partial t} + 2(\mathbf{U}\cdot\operatorname{grad})\xi_\varpi - 2\frac{U_\varphi\xi_\varphi}{\varpi} = \frac{\partial P}{\partial\varpi},
\end{equation}
\begin{equation}
\label{eq:12-245}
\frac{\partial\xi_\varphi}{\partial t} + 2(\mathbf{U}\cdot\operatorname{grad})\xi_\varphi + 2\frac{U_\varphi\xi_\varpi}{\varpi} = \frac{1}{\varpi}\frac{\partial P}{\partial\varphi};
\end{equation}
and
\begin{equation}
\label{eq:12-246}
\frac{\partial\xi_z}{\partial t} + 2(\mathbf{U}\cdot\operatorname{grad})\xi_z = \frac{\partial P}{\partial z},
\end{equation}
where
\begin{equation}
\label{eq:12-247}
\mathbf{U}\cdot\operatorname{grad} = U_\varpi\frac{\partial}{\partial\varpi} + \frac{U_\varphi}{\varpi}\frac{\partial}{\partial\varphi} + U_z\frac{\partial}{\partial z}.
\end{equation}
In addition we have the divergence condition
\begin{equation}
\label{eq:12-248}
\frac{\partial\xi_\varpi}{\partial\varpi} + \frac{\xi_\varpi}{\varpi} + \frac{1}{\varpi}\frac{\partial\xi_\varphi}{\partial\varphi} + \frac{\partial\xi_z}{\partial z} = 0.
\end{equation}

Analysing the disturbance into normal modes, we seek solutions whose $(t,\varphi,z)$-dependence is given by
\begin{equation}
\label{eq:12-249}
e^{i(\sigma t + m\varphi + kz)}.
\end{equation}
For $\mathbf{U}$ given by~\eqref{eq:12-243} and solutions having the chosen dependence on the coordinates,
\begin{equation}
\label{eq:12-250}
\frac{\partial}{\partial t} + 2\mathbf{U}\cdot\operatorname{grad} = i\left\{\sigma + 2U\left(\frac{m}{p} + k\right)\right\} = i\omega \quad \text{(say)},
\end{equation}
and the perturbation equations become
\begin{equation}
\label{eq:12-251}
i\omega\xi_\varpi - 2\frac{U}{p}\xi_\varphi = DP,
\end{equation}
\begin{equation}
\label{eq:12-252}
i\omega\xi_\varphi + 2\frac{U}{p}\xi_\varpi = \frac{im}{\varpi}P,
\end{equation}
\begin{equation}
\label{eq:12-253}
i\omega\xi_z = ikP,
\end{equation}
and
\begin{equation}
\label{eq:12-254}
D_*\xi_\varpi + \frac{im}{\varpi}\xi_\varphi + ik\xi_z = 0,
\end{equation}
where
\begin{equation}
\label{eq:12-255}
D = d/d\varpi \quad \text{and} \quad D_* = D + 1/\varpi.
\end{equation}
From equations~\eqref{eq:12-251} and \eqref{eq:12-252}, we find
\begin{equation}
\label{eq:12-256}
\left(\omega^2 - 4\frac{U^2}{p^2}\right)\xi_\varpi = -i\left(\omega DP + 2\frac{mU}{p}\frac{P}{\varpi}\right)
\end{equation}
and
\begin{equation}
\label{eq:12-257}
\left(\omega^2 - 4\frac{U^2}{p^2}\right)\xi_\varphi = +\left(2\frac{U}{p}DP + m\omega\frac{P}{\varpi}\right).
\end{equation}
Substituting for the components of $\boldsymbol\xi$ (in accordance with equations~\eqref{eq:12-253}, \eqref{eq:12-256}, and \eqref{eq:12-257}) in equation~\eqref{eq:12-254} and making use of the identity
\begin{equation}
\label{eq:12-258}
D_*D = D^2 + \frac{1}{\varpi}D,
\end{equation}
we obtain
\begin{equation}
\label{eq:12-259}
D_*DP + \frac{1}{\varpi}\left[\frac{4U^2/p^2}{\omega^2 - 4U^2/p^2}\right]^2 P = 0.
\end{equation}
The solution of equation~\eqref{eq:12-259} which is free of singularity at the origin is
\begin{equation}
\label{eq:12-260}
P = A J_m(\ell\varpi),
\end{equation}
where $A$ is a constant, $J_m$ is Bessel's function of order $m$, and
\begin{equation}
\label{eq:12-261}
\ell^2 = k^2\frac{\omega^2}{\omega^2 - 4U^2/p^2} - \frac{m^2}{\varpi^2}.
\end{equation}

If we impose on $P$ the same boundary condition as we did on the solution~\eqref{eq:12-259}), we shall obtain
\begin{equation}
\label{eq:12-262}
\omega = \frac{2U}{p}q,
\end{equation}
where $q$ has the meaning as in~\eqref{eq:12-261}. An equivalent form for $\sigma$ is
\begin{equation}
\label{eq:12-263}
\frac{\sigma R}{2U} = -\left(x+lm\right) \pm \frac{xl}{p}\sqrt{x^2+l^2_mq^2},
\end{equation}
which (since $q$ can be $\pm |q|$ for a given $x$ and $m$) is the same as the first of the two solutions included in~\eqref{eq:12-262}).

The preceding analysis shows how, in this case, the particular result~\eqref{eq:12-263}) is a consequence of the chosen boundary conditions. If we should suppose, instead, that the fluid is confined by a rigid wall at $\varpi = R$, the proper boundary condition would be that $u_\varpi$ vanishes at $\varpi = R$. Since (cf.\ equation~\eqref{eq:12-256})
\begin{equation}
\label{eq:12-264}
\left(\omega^2 - 4\frac{U^2}{p^2}\right)u_\varpi = i\sigma\left(\omega DP + 2\frac{mU}{p}\frac{P}{\varpi}\right),
\end{equation}
the vanishing of $u_\varpi$ at $\varpi = R$ will lead to the condition
\begin{equation}
\label{eq:12-265}
qJ'_m(y) + \frac{2mU}{pR}J_m(y) = 0,
\end{equation}
where
\begin{equation}
\label{eq:12-266}
y = \ell R = \sqrt{x^2 + p^2 q^2}\frac{R}{p}.
\end{equation}
Eliminating $2U/pR$ from equation~\eqref{eq:12-265}) by making use of equation~\eqref{eq:12-263}), we obtain
\begin{equation}
\label{eq:12-267}
yJ'_m(y) \pm m(1+p^2 x^2/y^2)J_m(y) = 0.
\end{equation}
We observe that equation~\eqref{eq:12-267}) is the same as the characteristic equation which governs the periods of oscillation of a rotating column of liquid (Chapter~VII, \S~68~(a)). The dispersion relations illustrated in Fig.~63 (p.~286) are, therefore, applicable to this problem with minor re-identifications.

If we denote by $y_{m,j}$ the $j$th zero of the transcendental equation~\eqref{eq:12-265}), the proper expression for $\sigma$ is
\begin{equation}
\label{eq:12-268}
\frac{\sigma R}{U} = -(x+lm) \pm 2\frac{xl}{}\sqrt{x^2+y^2_{m,j}},
\end{equation}
which apart from the replacement of $x_{m,j}$ by $y_{m,j}$, is the same as equation~\eqref{eq:12-263}). The main point is, however, that the results which follow from the use of the two different boundary conditions ($P = 0$ or $u_\varpi = 0$ at $\varpi = R$) are substantially the same: thus, for the modes $m = 0$, the change consists in the substitution of the zeros of $J_0(x)$ by the zeros of $J_1(x)$; and this is not very consequential.

\section{The effect of fluid motions on the stability of helical magnetic fields}\label{sec:114}
In \S~113 we have seen how the equipartition solution is stable under very general conditions. The essential characteristic of the solution is that in the stationary state both magnetic fields and fluid motions are present. It is, therefore, of interest to see, somewhat more generally than in the case of the equipartition solution, how fluid motions affect the stability of magnetic field configurations. A simple example of such more general configurations is when $\mathbf{u}$ and $\mathbf{h}$ are parallel to each other (as in the equipartition solution) but the constant of proportionality is such that equipartition of energy does not obtain.

When $\mathbf{u}$ and $\mathbf{h}$ are stationary and are parallel to each other, equation~\eqref{eq:12-218} governing the magnetic field is satisfied identically; but equation~\eqref{eq:12-217} cannot be satisfied so simply except when equipartition prevails. However, equation~\eqref{eq:12-217}) can be satisfied under special circumstances; thus
\begin{equation}
\label{eq:12-269}
\mathbf{u} = U\left(0, \frac{\varpi}{p}, 1\right) \quad \text{and} \quad \mathbf{h} = H\left(0, \frac{\varpi}{p}, 1\right)
\end{equation}
(where $U$ and $H$ are constants) satisfy equations~IX~(1)--(3) (with the terms in $\nu$ suppressed) provided
\begin{equation}
\label{eq:12-270}
\Pi = \text{constant} + \frac{1}{2}\left(U^2 - \frac{\mu H^2}{4\pi\rho}\right)\frac{\varpi^2}{p^2}.
\end{equation}

The stability of these general helical fields was first considered by Trehan.

Returning to equations~\eqref{eq:12-217} and \eqref{eq:12-218}, we obtain by addition and subtraction the symmetrical pair of equations
\begin{equation}
\label{eq:12-271}
\frac{\partial Y_j}{\partial t} + X_s\frac{\partial Y_j}{\partial x_s} + Y_s\frac{\partial X_j}{\partial x_s} = -\frac{\partial\Pi}{\partial x_j},
\end{equation}
and
\begin{equation}
\label{eq:12-272}
\frac{\partial X_j}{\partial t} + Y_s\frac{\partial X_j}{\partial x_s} + X_s\frac{\partial Y_j}{\partial x_s} = +\frac{\partial\Pi}{\partial x_j},
\end{equation}
where
\begin{equation}
\label{eq:12-273}
X_j = h_j - u_j \quad \text{and} \quad Y_j = h_j + u_j.
\end{equation}
Let equations~\eqref{eq:12-271} and \eqref{eq:12-272} allow a stationary solution of the form
\begin{equation}
\label{eq:12-274}
X_j = \alpha_0 U_j \quad \text{and} \quad Y_j = \alpha_1 U_j,
\end{equation}
where $\alpha_0$ and $\alpha_1$ are constants. (We shall presently identify $\mathbf{U}$ with a helical field of the type~\eqref{eq:12-269}.) Let $\xi_j$, $\zeta_j$, and $P$ represent a first order perturbation of $X_j$, $Y_j$, and $\Pi$, respectively. The equations governing the perturbations are
\begin{equation}
\label{eq:12-275}
\frac{\partial \xi_j}{\partial t} + \alpha_1 U_s\frac{\partial \xi_j}{\partial x_s} + \alpha_0 \zeta_s\frac{\partial U_j}{\partial x_s} = -\frac{\partial P}{\partial x_j},
\end{equation}
and
\begin{equation}
\label{eq:12-276}
\frac{\partial \zeta_j}{\partial t} + \alpha_0 U_s\frac{\partial \zeta_j}{\partial x_s} + \alpha_1 \xi_s\frac{\partial U_j}{\partial x_s} = +\frac{\partial P}{\partial x_j}.
\end{equation}
These equations are general. Now suppose that
\begin{equation}
\label{eq:12-277}
\mathbf{U} = U\left(0, \frac{\varpi}{p}, 1\right),
\end{equation}
where $U$ is a constant and the condition corresponding to~\eqref{eq:12-270} is satisfied.

Writing equations~\eqref{eq:12-275} and \eqref{eq:12-276} in cylindrical polar coordinates\footnote{This can be most readily accomplished by noting that in cylindrical polar coordinates, the vector $(A\cdot\text{grad})B$ has the components $(A\cdot\text{grad})B_\varpi - A_\varphi B_\varphi/\varpi$, $(A\cdot\text{grad})B_\varphi + A_\varphi B_\varpi/\varpi$, and $(A\cdot\text{grad})B_z$, where $(A\cdot\text{grad})$ is defined in the manner~\eqref{eq:12-247}. For $\mathbf{U}$ given by~\eqref{eq:12-277}) and for solutions having the chosen $(t,\varphi,z)$-dependence given by~\eqref{eq:12-249}, equations~\eqref{eq:12-275}--\eqref{eq:12-276}) are self-evident.} and seeking solutions whose $(t,\varphi,z)$-dependence is given by~\eqref{eq:12-249}, we obtain:
\begin{equation}
\label{eq:12-278}
i\omega_1\xi_\varpi - \alpha_0\frac{U}{p}\zeta_\varphi = DP,
\end{equation}
\begin{equation}
\label{eq:12-279}
i\omega_1\xi_\varphi + \alpha_0\frac{U}{p}\zeta_\varpi = \frac{im}{\varpi}P,
\end{equation}
\begin{equation}
\label{eq:12-280}
i\omega_1\xi_z = ikP,
\end{equation}
and
\begin{equation}
\label{eq:12-281}
i\omega_2\zeta_\varpi - \alpha_1\frac{U}{p}\xi_\varphi = -DP,
\end{equation}
where
\begin{equation}
\label{eq:12-282}
\omega_1 = \sigma + \alpha_1 U\left(\frac{m}{p}+k\right) \quad \text{and} \quad \omega_2 = \sigma + \alpha_0 U\left(\frac{m}{p}+k\right).
\end{equation}
Solving equations~\eqref{eq:12-278}--\eqref{eq:12-281} for $\xi_\varpi$ and $\xi_\varphi$ in terms of $DP$ and $P$, we readily find
\begin{equation}
\label{eq:12-283}
\left[\omega_1\omega_2 - \frac{U^2}{p^2}(\alpha_0 - \alpha_1)\omega_2\right]\xi_\varpi = i\omega_1\omega_2 DP + \frac{1}{2}(\alpha_0 - \alpha_1)\omega_2\frac{mU}{p}\frac{P}{\varpi},
\end{equation}
and
\begin{equation}
\label{eq:12-284}
\left[\omega_1\omega_2 - \frac{U^2}{p^2}(\alpha_0 - \alpha_1)\omega_2\right]\xi_\varphi = i\omega_1\omega_2\frac{m}{\varpi}P + DP - (\alpha_0 - \alpha_1)\frac{U}{p}\omega_2 DP.
\end{equation}
We also have
\begin{equation}
\label{eq:12-285}
\xi_z = \frac{k}{\omega_1}P.
\end{equation}
Now substituting for the components of $\boldsymbol\xi$ in accordance with the foregoing equations, in
\begin{equation}
\label{eq:12-286}
D_*\xi_\varpi + \frac{im}{\varpi}\xi_\varphi + ik\xi_z = 0
\end{equation}
the solenoidal character of $\boldsymbol\xi$, we obtain (on making use of the identity~\eqref{eq:12-258}))
\begin{equation}
\label{eq:12-287}
D_*DP + \left[\frac{k^2\omega_1\omega_2}{(\alpha_0-\alpha_1)\omega_2 U^2/p^2 - \omega_1\omega_2}\right]^2\frac{m^2}{\varpi^2}P = 0.
\end{equation}
The solution of this equation which is free of singularity at the origin is
\begin{equation}
\label{eq:12-288}
P = AJ_m(\ell\varpi),
\end{equation}
where $A$ is a constant, $J_m$ is Bessel's function of order $m$, and
\begin{equation}
\label{eq:12-289}
\ell^2 = k^2 + \frac{k^2\omega_1\omega_2}{(\alpha_0-\alpha_1)\frac{U^2}{p^2} - \omega_1\omega_2}.
\end{equation}

We shall suppose that the boundary condition\footnote{In general the boundary conditions may be more complex (see footnote on p.~556 and subsection~(b) below). The use of the boundary condition $P = 0$ at $\varpi = R$ is to be considered no more than as illustrative of the results which might follow.} require the vanishing of $P$ at $\varpi = R$. Then
\begin{equation}
\label{eq:12-290}
\ell R = x_{m,j},
\end{equation}
where $x_{m,j}$ denotes the $j$th root of $J_m(x)$.

The roots of this last equation are given by
\begin{equation}
\label{eq:12-291}
\sigma = -\frac{1}{2}(\alpha_0+\alpha_1)U\left(\frac{m}{p}+k\right) \pm \frac{1}{2}(\alpha_0-\alpha_1)U\sqrt{\left(\frac{m}{p}+k\right)^2 + \frac{x_{m,j}^2}{R^2}}.
\end{equation}
(Notice the double sign in the definition of $q$.)

It is now convenient to express the expression for $\sigma$ in terms of the fraction of the energy which is in the magnetic field. Since in the unperturbed state
\begin{equation}
\label{eq:12-292}
h_j - u_j = \alpha_0 U_j \quad \text{and} \quad h_j + u_j = \alpha_1 U_j,
\end{equation}
\begin{equation}
\label{eq:12-293}
f = \frac{\text{magnetic energy}}{\text{total energy}} = \frac{(\alpha_0+\alpha_1)^2}{(\alpha_0+\alpha_1)^2 + (\alpha_0-\alpha_1)^2}.
\end{equation}
Further, let (cf.\ equation~\eqref{eq:12-263}))
\begin{equation}
\label{eq:12-294}
l = R/p;
\end{equation}
and measure $\sigma$ in the unit
\begin{equation}
\label{eq:12-295}
\frac{U}{R}\sqrt{f(1+q^2)} \sec^{-1}.
\end{equation}
With these definitions, equation~\eqref{eq:12-291}) takes the form
\begin{equation}
\label{eq:12-296}
\sigma = -\frac{\sqrt{f}}{\sqrt{1+q^2}}[x + (q+m)l] \pm \frac{\sqrt{1-f}}{\sqrt{1+q^2}}\sqrt{[x+(q+m)l]^2 + x_{m,j}^2}.
\end{equation}

[Note that when $f = \frac{1}{2}$, this equation reduces to equation~\eqref{eq:12-263}) derived in \S~113~(a).]

Since for any given value of $x$ and $f$ we can associate with $q$ the two values $\pm|q|$, there are, in general, four distinct modes of oscillation corresponding to the four roots for $\sigma$ which follow from equation~\eqref{eq:12-296}).

From equation~\eqref{eq:12-296}) it is apparent that $\sigma$ cannot have a complex root if $f \leq \frac{1}{2}$. \emph{Accordingly, the system is stable if the energy in the magnetic field does not exceed the energy in the velocity field.} This result is due to Trehan.

Before we consider the full implications of equation~\eqref{eq:12-296}), we shall examine the two limiting cases, $f = 0$ and $f = 1$, when the energy is entirely of one kind or another.

\subsection*{(a) The case $f = 0$}

This is the case of a pure velocity field which can be obtained by setting $\alpha_0 = -\alpha_1$ and on the assumption that $\alpha_0$ is positive, equation~\eqref{eq:12-296}) gives
\begin{equation}
\label{eq:12-297}
\sigma = -[x+(2q+m)l]^{1/2}q;
\end{equation}
or
\begin{equation}
\label{eq:12-298}
\sigma = -[x+(2q+m)l] \quad \text{or} \quad -(x+m)l.
\end{equation}
However, a direct treatment of this problem leads only to the two modes represented by the first of the two solutions included in~\eqref{eq:12-298}). To see how this arises and to clarify, also, the possibility of other boundary conditions (besides the one we have used), we shall treat this case \emph{de novo}.

The perturbation equation in the case of a pure velocity field is
\begin{equation}
\label{eq:12-299}
\frac{\partial\xi_j}{\partial t} + U_s\frac{\partial\xi_j}{\partial x_s} + u_s\frac{\partial U_j}{\partial x_s} = -\frac{\partial P}{\partial x_j}.
\end{equation}

For solutions having the $(t,\varphi,z)$-dependence given by~\eqref{eq:12-249}) and for the same form of $\mathbf{U}$ as~\eqref{eq:12-277}), equation~\eqref{eq:12-299}) gives (see the footnote on p.~558)
\begin{equation}
\label{eq:12-300}
i\omega\xi_\varpi - \frac{U}{p}\xi_\varphi = -DP,
\end{equation}
\begin{equation}
\label{eq:12-301}
i\omega\xi_\varphi + \frac{U}{p}\xi_\varpi = -\frac{im}{\varpi}P,
\end{equation}
and
\begin{equation}
\label{eq:12-302}
i\omega\xi_z = -ikP,
\end{equation}
where
\begin{equation}
\label{eq:12-303}
\omega = \sigma + U\left(\frac{m}{p}+k\right).
\end{equation}
From the formal identity of the set of equations~\eqref{eq:12-300})--\eqref{eq:12-302}) with the set~\eqref{eq:12-251})--\eqref{eq:12-253}) and \eqref{eq:12-259}))
\begin{equation}
\label{eq:12-304}
P = AJ_m(\ell\varpi),
\end{equation}
where
\begin{equation}
\label{eq:12-305}
\ell^2 = k^2\frac{\omega^2}{\omega^2-U^2/p^2} - \frac{m^2}{\varpi^2}.
\end{equation}

If we impose on $P$ the same boundary condition as we did on the solution~\eqref{eq:12-259}), we shall obtain
\begin{equation}
\label{eq:12-306}
\omega = \frac{2U}{p}q,
\end{equation}
where $q$ has the meaning as in~\eqref{eq:12-261}). The corresponding solution for $\sigma$ is
\begin{equation}
\label{eq:12-307}
\sigma = -U\left(\frac{m}{p}+k\right) + \frac{2U}{p}q.
\end{equation}
An equivalent form of this equation is
\begin{equation}
\label{eq:12-308}
\frac{\sigma R}{U} = -(x+lm) + 2lq,
\end{equation}
which (since $q$ can be $\pm |q|$ for a given $x$ and $m$) is the same as the first of the two solutions included in~\eqref{eq:12-298}).

The preceding analysis shows how, in this case, the particular result~\eqref{eq:12-308}) is a consequence of the chosen boundary conditions. If we should suppose, instead, that the fluid is confined by a rigid wall at $\varpi = R$, the proper boundary condition would be that $u_\varpi$ vanishes at $\varpi = R$. Since (cf.\ equation~\eqref{eq:12-256}))
\begin{equation}
\label{eq:12-309}
\left(\omega^2 - \frac{U^2}{p^2}\right)u_\varpi = i\left(\omega DP + \frac{2mU}{pR}P\right),
\end{equation}
the vanishing of $u_\varpi$ at $\varpi = R$ will lead to the condition
\begin{equation}
\label{eq:12-310}
yJ'_m(y) + \frac{2mU}{p\omega}J_m(y) = 0,
\end{equation}
where
\begin{equation}
\label{eq:12-311}
y = \ell R = \sqrt{x^2+p^2q^2}\frac{1}{p}.
\end{equation}
Eliminating $2U/p$ from equation~\eqref{eq:12-310}) by making use of equation~\eqref{eq:12-306}), we obtain
\begin{equation}
\label{eq:12-312}
yJ'_m(y) \pm m(1+p^2x^2/y^2)^{1/2}J_m(y) = 0.
\end{equation}
We observe that equation~\eqref{eq:12-312}) is the same as the characteristic equation which governs the periods of oscillation of a rotating column of liquid (Chapter~VII, \S~68~(a)). The dispersion relations illustrated in Fig.~63 (p.~286) are, therefore, applicable to this problem with minor re-identifications.

If we denote by $y_{m,j}$ the $j$th zero of the transcendental equation~\eqref{eq:12-310}), the proper expression for $\sigma$ is
\begin{equation}
\label{eq:12-313}
\frac{\sigma R}{U} = -(x+lm) \pm 2\frac{xl}{\sqrt{x^2+y^2_{m,j}}},
\end{equation}
which apart from the replacement of $x_{m,j}$ by $y_{m,j}$, is the same as equation~\eqref{eq:12-308}). The main point is, however, that the results which follow from the use of the two different boundary conditions ($P = 0$ or $u_\varpi = 0$ at $\varpi = R$) are substantially the same: thus, for the modes $m = 0$, the change consists in the substitution of the zeros of $J_0(x)$ by the zeros of $J_1(x)$; and this is not very consequential.

\subsection*{(b) The case $f = 1$}

This is the case of a pure magnetic field which can be obtained by setting $\alpha_0 = \alpha_1$. Equation~\eqref{eq:12-296}) gives for this case
\begin{equation}
\label{eq:12-314}
\sigma = \pm\sqrt{f}\{x + (\pm|q|+m)l\} = \pm|q|l.
\end{equation}
From this equation for $\sigma$ it is evident that the condition for instability to occur is
\begin{equation}
\label{eq:12-315}
|x + (\pm|q|+m)l| < |q|l.
\end{equation}
From the definition of $q$ in equation~\eqref{eq:12-289}), it is apparent that $|q|$ is a monotone increasing function of $x$ and, moreover,
\begin{equation}
\label{eq:12-316}
|q| < 1,
\end{equation}
the value 1 being attained only for $x \to \infty$.

In view of~\eqref{eq:12-316}), the inequality~\eqref{eq:12-315}) can never be met for $m \geq 2$; therefore, \emph{the system is stable for all modes belonging to $m \geq 2$.} The other hand if $m$ should be negative, then with the choice of the negative sign in front of $|q|$ on the left-hand side of~\eqref{eq:12-315}), the inequality which must be satisfied for instability to occur is
\begin{equation}
\label{eq:12-317}
|x - (|q|+|m|)l| < |q|l,
\end{equation}
or
\begin{equation}
\label{eq:12-318}
(2|q|+|m|)l > x > (|q|+|m|)l.
\end{equation}
In the finite range of $x$ specified by~\eqref{eq:12-318}) the system is certainly unstable; therefore, \emph{there exist ranges of wave numbers for which the system is unstable for modes belonging to any negative $m$.}

It remains to consider the modes belonging to $m = 0$ and 1. When $m = 0$, the inequality~\eqref{eq:12-315}) becomes
\begin{equation}
\label{eq:12-319}
|x \pm |q|l| < |q|l.
\end{equation}
This inequality cannot be met with the positive sign on the left-hand side; therefore, it reduces to
\begin{equation}
\label{eq:12-320}
x < 2|q|l = \frac{2xl}{\sqrt{x^2+x^2_{0,1}}},
\end{equation}
or
\begin{equation}
\label{eq:12-321}
l > \frac{1}{2}x_{0,1}.
\end{equation}
A necessary condition that for some $x$ and $j$ a mode belonging to $m = 0$ is unstable is
\begin{equation}
\label{eq:12-322}
l = R/p > \frac{1}{2}x_{0,1} = 1.202.
\end{equation}

Finally, considering the case $m = 1$, we will have instability if
\begin{equation}
\label{eq:12-323}
|x + (1-|q|)l| < |q|l,
\end{equation}
a condition which cannot again be met with the positive sign in front of $|q|$. Therefore, the condition reduces to (in view of~\eqref{eq:12-316}))
\begin{equation}
\label{eq:12-324}
x + (1-|q|)l < |q|l
\end{equation}
or
\begin{equation}
\label{eq:12-325}
x < (2|q|-1)l.
\end{equation}

Inserting for $|q|$ its value from~\eqref{eq:12-289}), the condition becomes
\begin{equation}
\label{eq:12-326}
x < \left[\sqrt{x^2+x^2_{1,j}}-1\right]l.
\end{equation}
The right-hand side of~\eqref{eq:12-326}) is negative if
\begin{equation}
\label{eq:12-327}
x < x_{1,1}/\sqrt{3} < x_{1,1} = 2.212.
\end{equation}
Hence, for $x < 2.212$ the modes belonging to $m = -1$ are certainly stable; however, when this value is exceeded, instability is possible if $l$ is sufficiently small.

The main conclusion to be drawn from the foregoing discussion is that \emph{the helical magnetic field is unstable}. The remarkable fact is that fluid motions parallel to the field, if of sufficient intensity, can stabilize it.

\subsection*{(c) The general case}

We have seen that while equation~\eqref{eq:12-296}) predicts four distinct modes for a general value of $f$, there are only two modes for the particular values $f = 0$, $\frac{1}{2}$, and 1. The way in which the different modes coalesce and separate again as $f$ passes through these values is of interest.

When $f = 1$, we have the modes given by equation~\eqref{eq:12-314}). In this case, since we need not distinguish the $\pm$ sign which occurs as a factor in the equation, the two modes for $f = 1$ be distinguished by $1$ and $2$. When $f$ decreases, each of these two modes splits into two modes; we shall distinguish them by $1a$ and $1b$, and $2a$ and $2b$, respectively. In the case of a complex $\sigma$, the imaginary part of $\sigma$ decreases, while the real part increases as $f$ decreases. The imaginary part vanishes for some $f > \frac{1}{2}$.

\figurePlaceholder{130}{The dependence of the frequency of oscillation on the fraction of the total energy in the velocity field for the modes $m = 0$, and for $kR = 1$. The curves in (a), (b), and (c) are for $R/p = 1, \frac{1}{2}$, and $2$, respectively. The `hydromagnetic' and the `velocity' modes are shown by the solid and the dashed lines, respectively; while the unstable modes are distinguished by the dotted lines. For $R/p = \sqrt{2}$ and $2$, we have stability for $1-f > 0.236$ and $0.476$, respectively.}

\figurePlaceholder{131}{The dependence of the frequency of oscillation on the fraction of the total energy in the velocity field for the modes $m = -1$, and for $kR = 0.45$. The curves in (a), (b), and (c) are for $R/p = 1, \frac{1}{2}$, and $2$, respectively. The `hydromagnetic' and the `velocity' modes are shown by the solid and the dashed lines; respectively; while the unstable modes are distinguished by the dotted lines. For $R/p = 1$ and $\sqrt{2}$, we have stability for $1-f > 0.476$ and $0.228$, respectively.}

now all these modes are necessarily real. As $f$ tends to zero, the modes $1a$ and $2a$ coalesce while the remaining two modes tend to finite limits; and we obtain the three modes included in~\eqref{eq:12-298}). However, as we have seen, one of the modes included in~\eqref{eq:12-298}), namely, $\sigma = -(x+ml)$, is spurious; and one can, indeed, verify that in this limit the proper function belonging to this mode vanishes identically.

The dependence on $f$ of the four modes which follow from equation~\eqref{eq:12-296}) is shown in Figs.~130 (for the case $m = 0$) and 131 (for the case $m = -1$).

\section{The stability of the pinch}\label{sec:115}
The simplest of the so-called pinch configurations consists of a cylindrical column of fluid (the `plasma') inside of which a uniform axial magnetic field is present, while outside there is a similar axial field together with a transverse $\varphi$-field falling off as $\varpi^{-1}$. In the usual arrangements, the column of fluid is circled by a concentric conducting wall (see Fig.~132). Configurations of this kind are achieved in the laboratory by sending a high current through a gas by means of a discharge. The axial current produces a transverse magnetic field which `pinches' the column of gas into a configuration which we have idealized in the description; and the axial fields of different strengths inside and outside the column of fluid, which we have postulated, are `programmed' in the experimental arrangements. The stability of such pinch configurations and their achievement in stable form have acquired considerable practical significance in recent years.

Let $R_1$ be the radius of the fluid column and $R_2$ that of the encircling conducting wall. Let $H_i$ and $H_e$ denote the strengths of the axial magnetic fields inside and outside the column; and let the transverse $\varphi$-field be
\begin{equation}
\label{eq:12-328}
H_\varphi = H_\varphi R_1/\varpi.
\end{equation}
The fact that the magnetic field is discontinuous at $\varpi = R_1$ means that there is a current sheet of strength
\begin{equation}
\label{eq:12-329}
J_z = \frac{1}{4\pi}[(H_i - H_e)\mathbf{1}_\varphi + H_\varphi\mathbf{1}_z]
\end{equation}
on the surface. (Such current sheets occur when the fluid is assumed to be of infinite electrical conductivity, as in the present instance.) Furthermore, in the stationary state, the continuity of the normal stress across the surface of the fluid requires that the constant pressure $p_0$ inside the fluid is given by
\begin{equation}
\label{eq:12-330}
p_0 = \frac{\mu}{8\pi}(H_\varphi^2 + H_e^2 - H_i^2).
\end{equation}
It will be convenient to express $H_i$ and $H_e$ in terms of $H_\varphi$. Let
\begin{equation}
\label{eq:12-331}
H_i = \beta_1 H_\varphi \quad \text{and} \quad H_e = \beta_2 H_\varphi.
\end{equation}
An inequality which must obtain in virtue of~\eqref{eq:12-330}) is
\begin{equation}
\label{eq:12-332}
1 + \beta_2^2 \geq \beta_1^2.
\end{equation}

\figurePlaceholder{132}{The pinch configuration.}

The equations governing small perturbations of the fluid column are the same as equations~\eqref{eq:12-100})--\eqref{eq:12-103}). However, $\Pi$ has now the meaning
\begin{equation}
\label{eq:12-333}
\Pi = \frac{\delta p}{\rho} + \frac{\mu H h_z^{(\mathrm{in})}}{4\pi\rho},
\end{equation}
where we have distinguished $\mathbf{h}$ by $\mathbf{h}^{(\mathrm{in})}$ to denote that it signifies the perturbation in the field in the interior of the fluid. (We shall similarly distinguish the perturbation in the field in the exterior of the fluid by $\mathbf{h}^{(\mathrm{ext})}$.)

By considering a normal mode of disturbance by specifying that the deformed surface be governed by the equation
\begin{equation}
\label{eq:12-334}
\varpi = R_1 + \epsilon_0 e^{i(kz+m\varphi)+\sigma t},
\end{equation}
we seek solutions of the perturbation equations whose $(t,\varphi,z)$-dependence is given by
\begin{equation}
\label{eq:12-335}
e^{i(kz+m\varphi)+\sigma t}.
\end{equation}
The corresponding solution for $\Pi$ is
\begin{equation}
\label{eq:12-336}
\Pi = \epsilon_0\Pi_0 I_m(k\varpi)e^{ikz+im\varphi+\sigma t}.
\end{equation}
The integral of equation~\eqref{eq:12-101}) when $\mathbf{H}$ is axial is
\begin{equation}
\label{eq:12-337}
\mathbf{h}^{(\mathrm{in})} = \frac{ikH_i}{\sigma}\mathbf{u}.
\end{equation}
The equation of motion for $\mathbf{u}$ then gives
\begin{equation}
\label{eq:12-338}
\sigma\left(1 + \frac{k^2\mu H_i^2}{4\pi\rho\sigma^2}\right)\mathbf{u} = -\operatorname{grad}\Pi.
\end{equation}
In particular,
\begin{equation}
\label{eq:12-339}
u_\varpi = -\frac{\sigma}{\sigma^2+\Omega_{a,i}^2}\frac{\Pi_0}{k}kI'_m(k\varpi)e^{ikz+im\varphi+\sigma t},
\end{equation}
where $\Omega_{a,i}$ is the Alfv\'en frequency appropriate to $H_i$.

The perturbation in the magnetic field outside of the fluid is a vacuum field, and accordingly, under the present circumstances can be derived from a potential in the manner
\begin{equation}
\label{eq:12-340}
\mathbf{h}^{(\mathrm{ext})} = \operatorname{grad}[C_1 I_m(k\varpi) + C_2 K_m(k\varpi)]e^{ikz+im\varphi+\sigma t},
\end{equation}
where $C_1$ and $C_2$ are constants. For the components of $\mathbf{h}^{(\mathrm{ext})}$ we have
\begin{equation}
\label{eq:12-341}
h_\varpi^{(\mathrm{ext})} = k[C_1 I'_m(k\varpi) + C_2 K'_m(k\varpi)]e^{ikz+im\varphi+\sigma t},
\end{equation}
\begin{equation}
\label{eq:12-342}
h_\varphi^{(\mathrm{ext})} = \frac{im}{\varpi}[C_1 I_m(k\varpi) + C_2 K_m(k\varpi)]e^{ikz+im\varphi+\sigma t},
\end{equation}
\begin{equation}
\label{eq:12-343}
h_z^{(\mathrm{ext})} = ik[C_1 I_m(k\varpi) + C_2 K_m(k\varpi)]e^{ikz+im\varphi+\sigma t}.
\end{equation}

The solutions represented by equations~\eqref{eq:12-336}), \eqref{eq:12-337}), \eqref{eq:12-339}), and \eqref{eq:12-341})--\eqref{eq:12-343}) must satisfy a number of boundary conditions. These are: (1) the velocity $u_\varpi$ at $\varpi = R_1$ must be compatible with the assumed form of the deformed surface; (2) the radial component $h_\varpi^{(\mathrm{ext})}$ must vanish on the conducting wall at $\varpi = R_2$; (3) the component of the magnetic field normal to the deformed surface must be continuous; and (4) the normal component of the total stress tensor must be continuous across the deformed surface.

The first of the foregoing conditions requires
\begin{equation}
\label{eq:12-344}
\epsilon_0\sigma = -\frac{\epsilon_0\Pi_0}{\sigma^2+\Omega_{a,i}^2}xI'_m(x_1)e^{ikz+im\varphi+\sigma t} \quad = u_\varpi(R_1),
\end{equation}
or, using equation~\eqref{eq:12-339}),
\begin{equation}
\label{eq:12-345}
\sigma^2 + \Omega_{a,i}^2 = -\frac{\Pi_0}{\epsilon_0}x_1 I'_m(x_1),
\end{equation}
where
\begin{equation}
\label{eq:12-346}
x_1 = kR_1 \quad \text{and} \quad x_2 = kR_2.
\end{equation}
From equations~\eqref{eq:12-337}) and \eqref{eq:12-341}) the relation
\begin{equation}
\label{eq:12-347}
h_\varpi^{(\mathrm{in})}(R_1) = ik_\epsilon H_i e^{ikz+im\varphi+\sigma t} = h_\varpi^{(\mathrm{ext})}(R_1).
\end{equation}
The vanishing of $h_\varpi^{(\mathrm{ext})}$ at $\varpi = R_2$ gives (cf.\ equation~\eqref{eq:12-341}))
\begin{equation}
\label{eq:12-348}
C_1 = -C_2 \frac{K'_m(x_2)}{I'_m(x_2)}.
\end{equation}
Since the tangential components of the magnetic field at the boundary of the fluid, in the unperturbed state, are not continuous, we must allow for this fact in applying the third of the enumerated boundary conditions. When discontinuities in the tangential magnetic field are present in the stationary state, the boundary condition which must be applied in the perturbed state is\footnote{...}
\begin{equation}
\label{eq:12-349}
\delta N \cdot \text{jump}[\mathbf{H} + \mathbf{h}] = 0,
\end{equation}
where $\delta(f)$ signifies the jump which a quantity $f$ experiences at the boundary, $\mathbf{N}^{(0)}$ is the unit outward normal to the unperturbed boundary, $\delta N$ the change in $\mathbf{N}^{(0)}$ consequent to the perturbation, $\mathbf{H}^{(0)}$ the unperturbed magnetic field, and $\mathbf{h}$ the perturbation in $\mathbf{H}$. In the case under consideration
\begin{equation}
\label{eq:12-350}
\mathbf{N}^{(0)} = \mathbf{1}_\varpi; \quad \delta N = -ik\epsilon_0 e^{ikz+im\varphi+\sigma t}\mathbf{1}_z + i\frac{m}{R_1}\epsilon_0 e^{ikz+im\varphi+\sigma t}\mathbf{1}_\varphi,
\end{equation}
and
\begin{equation}
\label{eq:12-351}
\Delta(\mathbf{H}^{(0)}) = H_e \mathbf{1}_z + (H_e - H_i)\mathbf{1}_\varphi.
\end{equation}
The application of the condition~\eqref{eq:12-349}) therefore gives
\begin{equation}
\label{eq:12-352}
[h_\varphi^{(\mathrm{ext})} - h_\varphi^{(\mathrm{in})}]_{R_1} = \frac{im}{R_1}\epsilon_0(H_e - H_i)e^{ikz+im\varphi+\sigma t} = 0.
\end{equation}

Making use of equations~\eqref{eq:12-341}), \eqref{eq:12-343}), and \eqref{eq:12-348}), we obtain after simplifications
\begin{equation}
\label{eq:12-353}
\Pi_0 I_m(x_1) \left\{1 + \frac{x_1^2\Omega_{a,i}^2}{\sigma^2+\Omega_{a,i}^2}\frac{K'_m(x_1)I'_m(x_2)-I'_m(x_1)K'_m(x_2)}{K_m(x_1)I'_m(x_2)-I_m(x_1)K'_m(x_2)} \cdot \frac{I_m(x_1)}{I'_m(x_1)}\right\} = 0.
\end{equation}

Now eliminating $\Pi_0$ between equations~\eqref{eq:12-345}) and \eqref{eq:12-353}), we finally obtain the dispersion relation
\begin{equation}
\label{eq:12-354}
(\beta_1^2+1) x_1 + (\beta_2 x_1 \pm m)^2 \frac{K'_m(qx_1)I_m(x_1)-I_m(qx_1)K'_m(x_1)}{K_m(qx_1)I'_m(x_1)-I_m(qx_1)K'_m(x_1)} < \beta_2^2\frac{x_1^2 I_m(x_1)}{I'_m(x_1)},
\end{equation}
where $\Omega_{a,e}$ is the Alfv\'en frequency defined in terms of $H_e$.

\subsection*{(a) On stable pinch configurations}

The dispersion relation~\eqref{eq:12-354}) involves the three parameters
\begin{equation}
\label{eq:12-355}
\beta_1 = H_i/H_\varphi, \quad \beta_2 = H_e/H_\varphi, \quad \text{and} \quad q = x_2/x_1 = R_2/R_1.
\end{equation}
Since $\beta_1$ occurs as $\beta_1^2$ in~\eqref{eq:12-354}), the sign is clearly immaterial; and we shall suppose that it is positive. On the other hand, the sign of $\beta_2$ does matter; however, since a reversal of the sign of $\beta_2$ is equivalent to a reversal of the sign of $m$ and since both signs of $m$ are, moreover, allowed, there is again no loss of generality in supposing that $\beta_2$ is also positive. The parameters $\beta_1$ and $\beta_2$ are not unrestricted in view of~\eqref{eq:12-332}); also it is apparent that $q \geq 1$. The principal question of interest in connexion with the present problem is therefore: in what parts of the space
\begin{equation}
\label{eq:12-356}
q \geq 1, \quad \beta_1 \geq 0, \quad \beta_2 \geq 0, \quad \text{and} \quad 1+\beta_2^2 \geq \beta_1^2
\end{equation}
\emph{of the parameters do we have stability?}

A necessary and sufficient condition for stability can be given at once. From equation~\eqref{eq:12-354}) it is apparent that the principle of the exchange of stabilities is valid for the present problem and that \emph{a necessary and sufficient condition for stability is}
\begin{equation}
\label{eq:12-357}
x_1 + (\beta_2 x_1 \pm m)^2 \frac{K'_m(qx_1)I_m(x_1) - I_m(qx_1)K'_m(x_1)}{K_m(qx_1)I'_m(x_1) - I_m(qx_1)K'_m(x_1)} < \beta_2^2 \frac{x_1^2 I_m(x_1)}{I'_m(x_1)},
\end{equation}
\emph{for all $x_1$ $(> 0)$ and integral $m$ $(\geq 0)$}; for equation~\eqref{eq:12-354}), cannot, then, allow a solution with $\sigma^2 > 0$.

The states of marginal stability can be derived by considering~\eqref{eq:12-357}) with the equality sign in place of the inequality sign.

Two conclusions follow directly from~\eqref{eq:12-357}). The first, and from the practical standpoint of achieving stable pinch configurations an important, conclusion is that \emph{stability is impossible if $\beta_1 = 0$}. For when $\beta_1 x_1 = m$, the left-hand side of the inequality~\eqref{eq:12-357}) (with the negative sign for $m$ chosen) is positive, which contradicts~\eqref{eq:12-357}). The instability, in the absence of an internal $z$-field when $\beta_2 x_1 + m = 0$, occurs when the lines of force of the perturbed field are helices congruent to the one described by the unperturbed external field on the surface $\varpi = R_1$.

It is also clear that, other things being equal, stability for a negative $m$ implies stability for the corresponding positive $m$: for the additional term on the left-hand side of~\eqref{eq:12-357}) on passing from $m$ to $-m$ is
\begin{equation}
\label{eq:12-358}
-4|m|\beta_2 x_1 \frac{K'_m(qx_1)I_m(x_1)-I_m(qx_1)K'_m(x_1)}{K_m(qx_1)I'_m(x_1)-I_m(qx_1)K'_m(x_1)}
\end{equation}
and this is negative.

There is one further elementary result which we can derive. This concerns the case $m = 0$. When $m = 0$, the inequality~\eqref{eq:12-357}) becomes
\begin{equation}
\label{eq:12-359}
1-\beta_1^2\frac{I_0(x_1)K_0(x_1)+K_0(qx_1)I_0(x_1)}{I_0(x_1)K_0(qx_1)-K_0(x_1)I_0(qx_1)} < \beta_2^2\frac{x_1^2 I_0(x_1)}{I'_0(x_1)}.
\end{equation}
From the known properties of the Bessel functions, it can be shown that if the inequality~\eqref{eq:12-359}) holds for some value of $x_1$, then it holds for all larger values of $x_1$. Therefore, in this case, a necessary and sufficient condition that~\eqref{eq:12-359}) hold for all $x_1 > 0$ can be obtained by considering the limiting form of the inequality for $x_1 \to 0$. We thus obtain
\begin{equation}
\label{eq:12-360}
\frac{1}{2} - \frac{\beta_1^2}{q^2-1} < \beta_2^2.
\end{equation}

The foregoing is a necessary and sufficient condition that the modes belonging to $m = 0$ are stable.

An analytical discussion of the equation governing marginal stability for values of $m$ besides zero is possible. However, since in this particular instance exhaustive numerical results are available, we shall content ourselves by giving an account of them. The particular results which we shall describe and illustrate are those of Tayler.

When there is no encircling wall, $q \to \infty$, and the equation for marginal stability reduces to
\begin{equation}
\label{eq:12-361}
x_1 + (\beta_2 x_1 \pm m)^2 \frac{K_m(x_1)}{I'_m(x_1)} = \beta_2^2 \frac{x_1^2 I_m(x_1)}{I'_m(x_1)}.
\end{equation}
From this equation it appears that by letting $\beta_2$ exceed a certain determinate value (depending on $\beta_1$, except, as we shall see, in the case $m = 0$), one can achieve stability for all modes except those belonging to $m = -1$.

\begin{table}[ht]
\centering
\caption{Critical values of $\beta_1$ for the stability of the modes belonging to different $m$'s in case there is no encircling conducting wall}
\label{tab:LXIV}
\begin{tabular}{ccccccc}
\hline
& $\beta_2$ & 0.0 & 0.5 & 1.0 & 1.5 & 2.0 \\
\hline
$m$ & & & & & & \\
0 & & 0.707 & 0.707 & 0.707 & 0.707 & 0.707 \\
2 & & 0.539 & 0.696 & 1.121 & 1.521 & 2.023 \\
3 & & 0.169 & 0.492 & 0.801 & 1.130 & 1.468 \\
4 & & 0.126 & 0.499 & 1.122 & 1.592 & 2.073 \\
\hline
\end{tabular}
\end{table}

Thus, the modes belonging to $m = 0$ can be stabilized by the choice (see equation~\eqref{eq:12-360}))
\begin{equation}
\label{eq:12-362}
\beta_1^2 > 0.5
\end{equation}
independently of $\beta_2$. There are similar results for other values of $-|m|$ (depending, however, on $\beta_2$). These limits are given in Table~LXIV and are exhibited in Fig.~133. In Fig.~133 the curve $\beta_1^2 = 1+\beta_2^2$ and the loci of the marginally stable configurations for $m = 0$ and $m = -2$ bound a region (shown shaded in the figure) in which all modes belonging to $m \neq -1$ are stable. In the same shaded region, disturbances belonging to $m = -1$ and wave numbers not exceeding certain determinate values are unstable.

\figurePlaceholder{133}{Stability diagram for pinch configurations when there are no conducting walls. The modes belonging to $\pm m$ are stable when the representative point in the plane is above the curve labelled by $m = -|m|$ and below the locus $1+\beta_2^2 = \beta_1^2$. Complete stability for all modes except those belonging to $m = -1$ is obtained for values of $\beta_2$ and $\beta_1$ in the shaded region.}

\begin{table}[ht]
\centering
\caption{Wave numbers below which the modes belonging to $m = -1$ are unstable (no encircling conducting wall)}
\label{tab:LXV}
\begin{tabular}{ccccccc}
\hline
$\beta_2$ & 0.0 & 0.500 & 1.000 & 1.500 & 2.000 \\
\hline
$\beta_1$ & 1.000 & 1.118 & 1.414 & 1.803 & 2.236 \\
$x_1$ & 0.450 & 0.869 & 0.736 & 0.571 & 0.455 \\
$\beta_1$ & 0.707 & 0.707 & 1.122 & 1.592 & 2.073 \\
$x_1$ & 1.025 & 1.886 & 0.991 & 0.661 & 0.497 \\
\hline
\end{tabular}
\end{table}

If an encircling wall is allowed, it is found that the modes belonging to $m = -1$ can also be stabilized if $q$ $(> 1)$ is less than a certain critical value. Table~LXVI gives the maximum values of $q$ which will ensure stability for given values of $\beta_2$ when $\beta_1$ has the maximum permissible value, namely, $\sqrt{1+\beta_2^2}$. We observe that for $\beta_2 = 0$, the maximum value of $q$ compatible with stability is 5; when $\beta_2$ increases, the allowed range of $q$ decreases very rapidly. For a given $\beta_2$, when $\beta_1$ decreases

\begin{table}[ht]
\centering
\caption{Critical wall radii for the stability of the modes belonging to $m = -1$}
\label{tab:LXVI}
\begin{tabular}{cccccc}
\hline
$\beta_2$ & 0.0 & 0.100 & 0.200 & 0.500 & 1.000 \\
$\beta_1$ & 1.000 & 1.005 & 1.020 & 1.118 & 1.414 \\
\hline
\end{tabular}
\end{table}

below $\sqrt{1+\beta_2^2}$, the maximum allowed value of $q$, for stability, also decreases. For $q = 1$, the minimum values of $\beta_1$ required for the stability of the modes belonging to $m = -1$ when $q = 1$, are given in Table~LXVII. The results for this and other values of $q$ are included in Fig.~134; this figure contains the answer to the principal question raised at the beginning of this section.

\begin{table}[ht]
\centering
\caption{Minimum values of $\beta_1$ for the stability of the modes belonging to $m = -1$ when $q = 1$}
\label{tab:LXVII}
\begin{tabular}{ccccccc}
\hline
$\beta_2$ & 0.0 & 0.500 & 0.500 & 1.000 & 1.500 & 2.000 \\
$\beta_1$ & 0.000 & 0.495 & 0.583 & 1.114 & 1.580 & 0.061 \\
\hline
\end{tabular}
\end{table}

\figurePlaceholder{134}{Stability diagram for pinch configurations when conducting walls are present at a radius $R_2 = qR_1$. For a given value of $q$, the modes belonging to $m = 0$ are stable when the representative point in the plane is above the curve $A$; and the modes belonging to $m = -1$ are stable when the representative point is above the curve $B$. Complete stability for a given $q$ is obtained when the representative point lies in the region appropriate for the stability of the modes belonging to both $m = 0$ and $-1$.}

\section*{BIBLIOGRAPHICAL NOTES}\label{sec:12-bib}

Interest in the stability of liquid jets was stimulated by some beautiful experiments of Savart and their interpretation by Plateau in terms of capillarity. And Rayleigh laid the foundations for the theoretical treatment of these and similar problems in:

\begin{enumerate}
\item Lord \textsc{Rayleigh}, `On the instability of jets', \emph{Scientific Papers}, i, 361--71, Cambridge, England, 1899.
\end{enumerate}
In this paper, Rayleigh developed the important concept of the mode of maximum instability. An account of all this early work will be found in:
\begin{enumerate}
\setcounter{enumi}{1}
\item Lord \textsc{Rayleigh}, \emph{Theory of Sound}, 357--64, Dover Publications, New York, 1945.
\end{enumerate}
A delightful lecture including some beautiful illustrations is:
\begin{enumerate}
\setcounter{enumi}{2}
\item Lord \textsc{Rayleigh}, `Some applications of photography', \emph{Scientific Papers}, iii, 441--51, Cambridge, England, 1902.
\end{enumerate}
The recommendation of terms `varicose' and `varicosity' is contained in this lecture.

The interest in the stability of cylindrical systems has been revived in recent years by certain astronomical and physical problems: the stability of the spiral arms of galaxies and the stability of the so-called pinch configurations.

\S~108. The gravitational instability of an infinite cylinder (in a larger context) is considered in:
\begin{enumerate}
\setcounter{enumi}{3}
\item S. \textsc{Chandrasekhar} and E. \textsc{Fermi}, `Problems of gravitational stability in the presence of a magnetic field', \emph{Astrophys.\ J.}\ \textbf{118}, 116--41 (1953).
\end{enumerate}
See also:
\begin{enumerate}
\setcounter{enumi}{4}
\item S. \textsc{Chandrasekhar}, `Problems of stability in hydrodynamics and hydromagnetics', \emph{Monthly Notices Roy.\ Astron.\ Soc.\ London}, \textbf{113}, 667--78 (1953).
\end{enumerate}

\S~109. The role of viscosity in capillary instability was considered by Rayleigh (see reference~11 below and the comments on it); the treatment in the text is an extension to gravitational instability.

\S~110. The stabilizing effect of an axial magnetic field on the gravitational instability of an infinite cylinder is considered in reference~4. To illustrate the nature of the effect under the simplest conditions, the case considered in which the field was entirely confined within the cylinder and none existed outside it. Extensions of the treatment in reference~4 to other distributions of the field (including the one considered in the text) are contained in:
\begin{enumerate}
\setcounter{enumi}{5}
\item F. C. \textsc{Auluck} and D. S. \textsc{Kothari}, `The influence of a magnetic field on the longitudinal stability of a gravitating cylinder', \emph{Z.\ f.\ Astrophysik}, \textbf{42}, 101--13 (1957).

\item S. K. \textsc{Trehan}, `The stability of an infinitely long cylinder with a prevalent force-free magnetic field', \emph{Astrophys.\ J.}\ \textbf{127}, 436--45 (1958).

\item R. \textsc{Simon}, `The hydromagnetic oscillations of an incompressible cylinder', ibid.\ \textbf{128}, 375--83 (1958).

\item \rule{1cm}{0.4pt} `Les oscillations hydromagn\'etiques d'un cylindre incompressible', \emph{Ann.\ d'Astrophys.}\ \textbf{22}, 712--26 (1959).
\end{enumerate}

\S~111. As we have stated, the classical investigations on capillary instability are those of Savart, Plateau, and Rayleigh.

\S~111~(b). The capillary instability of a hollow jet is treated in:
\begin{enumerate}
\setcounter{enumi}{9}
\item Lord \textsc{Rayleigh}, `On the instability of cylindrical fluid surfaces', \emph{Scientific Papers}, iii, 594--6, Cambridge, England, 1902.
\end{enumerate}

\S~111~(c). The effect of viscosity was also treated by Rayleigh:
\begin{enumerate}
\setcounter{enumi}{10}
\item Lord \textsc{Rayleigh}, `On the instability of a cylinder of viscous liquid under capillary force', \emph{Scientific Papers}, iii, 585--93, Cambridge, England, 1902.
\end{enumerate}
However, Rayleigh restricted himself to the case when the effect of viscosity was paramount. The general case does not seem to have been considered before.

\S~112. The problem considered in this section does not seem to have been investigated, though some beautiful experiments bearing on this topic have been reported by:
\begin{enumerate}
\setcounter{enumi}{11}
\item A. \textsc{Dattner}, B. \textsc{Lehnert}, and S. \textsc{Lundquist}, `Liquid conductor model of instabilities in a pinched discharge', \emph{Second UN Geneva Conference}, 355--7, Pergamon Press, London, 1959.
\end{enumerate}

\S~113. The stability of the simplest solution of the equations of hydromagnetics is proved in:
\begin{enumerate}
\setcounter{enumi}{12}
\item S. \textsc{Chandrasekhar}, `On the stability of the simplest solution of the equations of hydromagnetics', \emph{Proc.\ Nat.\ Acad.\ Sci.}\ \textbf{42}, 273--6 (1956).
\end{enumerate}
The example considered in \S~(a) is due to:
\begin{enumerate}
\setcounter{enumi}{13}
\item S. K. \textsc{Trehan}, `The hydromagnetic oscillations of twisted magnetic fields, I', \emph{Astrophys.\ J.}\ \textbf{127}, 446--53 (1958).
\end{enumerate}
See also:
\begin{enumerate}
\setcounter{enumi}{14}
\item \rule{1cm}{0.4pt} and W. H. \textsc{Reid}, `The hydromagnetic oscillations of twisted magnetic fields.\ II', ibid.\ 454--8 (1958).
\end{enumerate}

\S~114. The general problem treated in this section was first considered by:
\begin{enumerate}
\setcounter{enumi}{15}
\item S. K. \textsc{Trehan}, `The effect of fluid motions on the stability of twisted magnetic fields', ibid.\ \textbf{129}, 475--82 (1959).
\end{enumerate}
The case when no motions are present was treated by:
\begin{enumerate}
\setcounter{enumi}{16}
\item P. H. \textsc{Roberts}, `Twisted magnetic fields', ibid.\ \textbf{124}, 430--42 (1956).
\end{enumerate}
See also:
\begin{enumerate}
\setcounter{enumi}{17}
\item S. \textsc{Lundquist}, `On the stability of magneto-hydrostatic fields', \emph{Physical Review}, \textbf{83}, 307--11 (1951).

\item J. W. \textsc{Dungey} and R. E. \textsc{Loughhead}, `Twisted magnetic fields in conducting fluids', \emph{Australian J.\ Phys.}\ \textbf{7}, 5--13 (1954).
\end{enumerate}

\S~115. The literature on the stability of pinch configurations is an extensive one. The following is a very brief list of those papers which bear on the very special problem treated in the text:
\begin{enumerate}
\setcounter{enumi}{19}
\item M. \textsc{Kruskal} and J. L. \textsc{Tuck}, `The instability of a pinched fluid with a longitudinal magnetic field', \emph{Proc.\ Roy.\ Soc.\ (London)} A, \textbf{245}, 222--37 (1958).

\item M. \textsc{Rosenbluth}, `Stability of the pinch', \emph{Los Alamos Scientific Laboratory Report}, No.~2030, April 1956.

\item R. J. \textsc{Tayler}, `A note on hydromagnetic stability problems', \emph{Phil.\ Mag.}, Ser.~8, \textbf{2}, 33--36 (1957).

\item \rule{1cm}{0.4pt} `Hydromagnetic instabilities of an ideally conducting fluid', \emph{Proc.\ Phys.\ Soc.\ London}, B, \textbf{70}, 31--48 (1957).

\item \rule{1cm}{0.4pt} `The influence of an axial magnetic field on the stability of a constricted gas discharge', ibid.\ 1049--63 (1957).

\item V. D. \textsc{Shafranov}, `On the stability of a plasma column in the presence of a longitudinal magnetic field and a conduction casing', \emph{Plasma Physics and the Problem of Controlled Thermonuclear Reactions}, ii, 197--215, edited by M. A. Leontovich, Pergamon Press, 1959.

\item T. F. \textsc{Volkov}, `On stability of a plasma cylinder in an external magnetic field', ibid.\ 216--22, edited by M. A. Leontovich, Pergamon Press, 1959.

\item K. \textsc{Hain}, R. \textsc{L\"ust}, and A. \textsc{Schl\"uter}, `Zur Stabilit\"at eines Plasmas', \emph{Z.\ f.\ Naturforschung}, \textbf{12a}, 833--41 (1957).

\item \rule{1cm}{0.4pt} \rule{1cm}{0.4pt} `Zur Stabilit\"at zylindersymmetrischer Plasmakonfigurationen mit Volumenstr\"omen', ibid.\ \textbf{13a}, 936--40 (1958).

\item S. \textsc{Chandrasekhar}, A. N. \textsc{Kaufman}, and K. M. \textsc{Watson}, `The stability of the pinch', \emph{Proc.\ Roy.\ Soc.\ (London)} A, \textbf{245}, 435--55 (1958).
\end{enumerate}

The numerical results described and illustrated in the text are those of Tayler (reference~24).
